\documentclass[11pt,a4paper,onecolumn]{article}

\usepackage[UTF8,scheme=plain]{ctex}
\usepackage[margin=2.5cm]{geometry}
\usepackage{cite}
\usepackage{booktabs}
\usepackage{array}
\usepackage{tabularx}
\usepackage{multirow}
\usepackage{url}
\usepackage{xcolor}
\usepackage{amsmath}
\usepackage{pifont}
\usepackage{enumitem}
\usepackage{tikz}
\usetikzlibrary{shapes.geometric,arrows.meta,positioning,fit,backgrounds,calc,trees,matrix,decorations.pathreplacing,shadows}
\usepackage{hyperref}
\usepackage{caption}
\captionsetup{font=small,labelfont=bf,labelsep=colon,skip=6pt}

% Table styling
\usepackage{colortbl}
\definecolor{tabhead}{RGB}{52,73,94}
\definecolor{tabrow}{RGB}{236,240,241}
\definecolor{tabalt}{RGB}{248,249,250}
\newcommand{\tabheader}[1]{\textcolor{white}{\textbf{#1}}}
\newcommand{\hcolor}[2]{\cellcolor{#1}\textcolor{white}{\textbf{#2}}}
\renewcommand{\arraystretch}{1.3}

\hypersetup{
  colorlinks=true,
  linkcolor=blue,
  citecolor=blue,
  urlcolor=blue
}

\newcommand{\cmark}{\ding{51}}
\newcommand{\xmark}{\ding{55}}
\newcommand{\risk}[1]{\texttt{#1}}
\newcommand{\code}[1]{\texttt{#1}}

\title{\textbf{大模型辅助应用程序审计方法综述}\\[0.3em]
\large A Survey of Large Language Model Assisted Application Auditing}

\author{匿名作者\\
\small 代码审计与可信执行环境研究方向}

\date{}

\begin{document}
\maketitle

\begin{abstract}
大语言模型（Large Language Models, LLMs）已从通用语言生成工具演化为程序分析与安全审计中的语义组件。本文基于 200 余篇文献和 158 条精选引用，构建面向应用程序审计的结构化综述框架，将现有方法归纳为十类：直接 LLM 代码审查、漏洞检测与修复、LLM 与静态分析协同、检索增强与结构化表示增强、智能体/执行链审计、prompt 与安全感知微调、benchmark 与评估协议、模型完整性与后门防御、LLM 增强安全测试、以及 LLM 代码生成安全。

本文通过 5 张结构图和 10 张比较表系统呈现：（1）方法分类体系图展示十类方法的层次关系；（2）技术演进时间线标注 2013--2026 关键里程碑；（3）多阶段审计流水线给出推荐架构与量化目标；（4）Agent 攻防全景图映射 5 类攻击与 5 类防御；（5）Benchmark 覆盖热力图揭示模型后门维度的评测空白。

关键发现包括：（1）LLM 审查接受率仅 18--84\%，数据噪声是关键瓶颈；（2）IEEE S\&P 和 USENIX Security 研究证明 LLM 无法可靠识别安全漏洞，更大模型不优于小模型；（3）间接提示注入可使 agent 防御降低 30\% 以上效用；（4）越狱在统计意义上不可避免；（5）模型后门和 AI 生成代码扩展了审计攻击面。本文建立"结论--文献证据--benchmark 样本--评测工件"的追踪矩阵，结合本仓库 100 样本审计 benchmark 的迭代结果，提出高可靠审计应从单轮问答转向证据链审计，从单模型判断转向多阶段系统，从代码可疑性转向行为可疑性，并将不确定性和 LLM 不可用显式纳入 fail-closed 机制。
\end{abstract}

\noindent\textbf{关键词：}大语言模型，代码审查，漏洞检测，程序分析，静态分析，检索增强生成，智能体安全，越狱攻击，提示注入，应用程序审计，安全软件工程
\vspace{0.5em}

\section{引言}

\subsection{问题背景}
应用程序审计（application auditing）指对软件代码、配置、依赖、运行逻辑与外部交互行为进行安全性分析，以识别命令执行、数据外传、敏感信息泄露、持久化、权限越界、提示注入、工具滥用等风险。传统审计主要依赖人工阅读、规则引擎、污点分析、静态分析和符号/抽象执行等方法。这些方法在既知模式识别上较强，但在跨文件推理、语义恢复、工程上下文补全和自然语言证据整合方面存在限制。

LLM 的引入改变了审计流程。与只做模式匹配的规则系统不同，LLM 可以解释代码意图、调用链语义、框架惯例和潜在风险目的，从而在安全审计中扮演语义解释器、候选筛选器、风险裁决器或证据整合器等角色。问题在于，LLM 的"语义能力"并不自动转化为"安全可靠性"：模型可能对可疑模式过度敏感，导致误报；也可能在证据不足时给出错误但自信的放行判断，导致漏报。多项实证研究表明，LLM 在安全漏洞识别中的可靠性仍不足以替代专业审计工具~\cite{ullaha_vuln_reasoning_2024,paquette2024llmscannot,risse2024limits}：六个顶级 ML4VD 模型区分漏洞函数与其修复补丁的准确率低于随机猜测~\cite{risse2024limits}，更大模型反而不如小模型（LLaMA-3 70B F1=5.10\% vs Gemma 7B F1=57.98\%）~\cite{niu2025largetomammoth}，agent 平均 ASR 高达 84.30\% 且能力越强的模型越脆弱~\cite{zhang2024asb}。

因此，LLM 在应用程序审计中的最佳角色并非替代传统方法，而是嵌入多阶段审计体系：静态分析负责召回，结构化证据负责上下文压缩，LLM 负责语义裁决，策略层负责最终阻断或放行，人工复核处理边界样本。

图~\ref{fig:timeline} 展示了代码审计技术从传统方法到 LLM 时代的演进路线。

\begin{figure}[htbp]
\centering
\begin{tikzpicture}
% Color definitions
\definecolor{era1}{RGB}{70,130,180}    % Steel blue
\definecolor{era2}{RGB}{210,140,60}    % Warm orange
\definecolor{era3}{RGB}{180,60,60}     % Deep red

% Main timeline
\draw[line width=2pt,color=black!30] (0,0) -- (15,0);
\draw[->,>=latex,line width=2pt,color=black!50] (15,0) -- (15.5,0);

% Era backgrounds
\fill[era1!8,rounded corners=4pt] (0,-0.4) rectangle (5,3.8);
\fill[era2!8,rounded corners=4pt] (5,-0.4) rectangle (9,3.8);
\fill[era3!8,rounded corners=4pt] (9,-0.4) rectangle (15,3.8);

% Era labels
\node[font=\small\bfseries,text=era1] at (2.5,3.5) {传统方法};
\node[font=\small\bfseries,text=era2] at (7,3.5) {深度学习};
\node[font=\small\bfseries,text=era3] at (12,3.5) {LLM 时代};

% Year ticks
\foreach \x/\year in {0/2013, 2.5/2017, 5/2019, 7/2021, 9/2023, 11/2024, 13/2025, 15/2026} {
  \draw[line width=1pt,color=black!40] (\x,-0.15) -- (\x,0.15);
  \node[font=\scriptsize,text=black!60,below] at (\x,-0.25) {\year};
}

% Milestones - Era 1 (above line)
\fill[era1] (0,0) circle(3pt);
\node[font=\scriptsize,anchor=south west,text=era1!80!black] at (0.15,0.3) {MCR 基础~\cite{bacchelli_bird_2013}};

\fill[era1] (2.5,0) circle(3pt);
\node[font=\scriptsize,anchor=south west,text=era1!80!black] at (2.65,0.3) {VUDDY~\cite{kim2017vuddy}};

\node[font=\scriptsize,anchor=south west,text=era1!80!black] at (2.65,0.75) {VulDeePecker~\cite{li2018vuldeepecker}};

% Milestones - Era 2
\fill[era2] (5,0) circle(3pt);
\node[font=\scriptsize,anchor=south west,text=era2!80!black] at (5.15,0.3) {Devign~\cite{zhou2019devign}};

\fill[era2] (6,0) circle(3pt);
\node[font=\scriptsize,anchor=south west,text=era2!80!black] at (6.15,0.3) {CodeBERT~\cite{feng2020codebert}};

\fill[era2] (7,0) circle(3pt);
\node[font=\scriptsize,anchor=south west,text=era2!80!black] at (7.15,0.3) {LineVul~\cite{fu2021linevul}};

% Milestones - Era 3
\fill[era3] (9,0) circle(3pt);
\node[font=\scriptsize,anchor=south west,text=era3!80!black] at (9.15,0.3) {SecLLMHolmes~\cite{paquette2024llmscannot}};

\fill[era3] (10,0) circle(3pt);
\node[font=\scriptsize,anchor=south west,text=era3!80!black] at (10.15,0.3) {InjecAgent~\cite{zhan2024injecagent}};

\node[font=\scriptsize,anchor=south west,text=era3!80!black] at (10.15,0.75) {MasterKey~\cite{deng2024masterkey}};

\fill[era3] (11,0) circle(3pt);
\node[font=\scriptsize,anchor=south west,text=era3!80!black] at (11.15,0.3) {AgentDojo~\cite{debenedetti2024agentdojo}};

\fill[era3] (12,0) circle(3pt);
\node[font=\scriptsize,anchor=south west,text=era3!80!black] at (12.15,0.3) {LLMxCPG~\cite{llmxcpg_2025}};

\node[font=\scriptsize,anchor=south west,text=era3!80!black] at (12.15,0.75) {AppAtch~\cite{wen2025appatch}};

\fill[era3] (13,0) circle(3pt);
\node[font=\scriptsize,anchor=south west,text=era3!80!black] at (13.15,0.3) {StruQ~\cite{chen2025struq}};

\node[font=\scriptsize,anchor=south west,text=era3!80!black] at (13.15,0.75) {IsolateGPT~\cite{wu2025isolateGpt}};

% Trend arrow
\draw[->,>=latex,line width=1.5pt,color=black!20,decorate,decoration={brace,amplitude=5pt,mirror}] (0.5,2.5) -- (14.5,2.5);
\node[font=\footnotesize\itshape,text=black!40] at (7.5,2.2) {审计复杂度与攻击面持续扩大};

\end{tikzpicture}
\caption{代码审计技术演进时间线（2013--2026）}
\label{fig:timeline}
\end{figure}

\subsection{术语界定}
本文区分四个层次：（1）代码审查，面向函数、补丁、提交或 Pull Request 的局部判断；（2）漏洞检测，面向可利用缺陷的位置、类型和证据链识别；（3）程序分析，面向调用关系、数据流、控制流和语义解释的结构化理解；（4）应用程序审计，覆盖源码、配置、依赖、检索、工具调用、记忆与运行时行为链。本文关注第四层，并将前三层方法视为应用程序审计的组成部分。

\subsection{研究问题与贡献}
本文围绕四个研究问题展开：
\begin{itemize}[leftmargin=*]
  \item RQ1：LLM 辅助应用程序审计的主要方法家族有哪些？
  \item RQ2：不同方法家族依赖什么证据、解决什么问题、受什么约束？
  \item RQ3：在模型选型、prompt 设计、fine-tuning、结构化证据和 benchmark 评估方面，哪些实践更稳定？
  \item RQ4：当前文献揭示了哪些共性局限，未来研究方向是什么？
\end{itemize}

本文贡献如下：
\begin{enumerate}[leftmargin=*]
  \item 提出面向应用程序审计的七类方法分类，并新增越狱攻防专题；
  \item 从证据类型、输出粒度、部署方式和复现约束比较不同方法，覆盖 200 余篇文献；
  \item 总结 LLM 审计的常见失效模式与威胁有效性问题；
  \item 建立"结论--文献证据--benchmark 样本--评测工件"的追踪矩阵；
  \item 给出面向高可靠审计系统的研究议程与工程建议。
\end{enumerate}

\section{综述方法}

\subsection{文献来源与检索范围}
本文参考文献覆盖 200 余篇，主要来自 IEEE S\&P、USENIX Security、NDSS、NeurIPS、ACL、EMNLP、ICSE、ASE、ISSTA、ACM TOSEM、ACM CSUR、Springer 与 arXiv 预印本，重点覆盖 2023--2026 年发表的工作。检索关键词包括：code review、secure code review、vulnerability detection、vulnerability repair、software security、static analysis、program analysis、retrieval-augmented generation、agent security、prompt injection、jailbreak、malicious skills、benchmark、fair evaluation、secure-aware fine-tuning 与 prompt mixture。

\subsection{纳入、排除与证据等级}
纳入标准为：与 LLM、代码、安全审计、漏洞检测、程序分析、智能体安全或相关评估方法直接相关；提供方法、系统、综述或实证评估；能够为应用程序审计提供方法论启示。排除标准为：仅讨论通用代码生成而无安全分析内容；与软件安全关联过弱；仅有标题或摘要相关性但缺少方法信息。

为避免把标题相关误认为证据充分，本文将参考文献按证据强度分为三层：A 级为综述或系统综述，用于归纳领域结构与研究空白~\cite{sheng2026llmsecurity,chen2026lmcode,kaniewski2026detecting,zhou2025vulndetrepair,sheng2025cybersecurity_slr,davila_nunes_2021_slr_taxonomy,code_security_slr_2024,lm_code_security_slr_2024,llm_security_survey_2025,slr_vuln_llm_2025}；B 级为实证或对比研究，用于判断实际可用性与稳定性~\cite{ullaha_vuln_reasoning_2024,paquette2024llmscannot,risse2024limits,niu2025largetomammoth,charoenwet_2024_sast_secure_review}；C 级为系统、工具或新方法论文，用于观察工程趋势与技术路线~\cite{llmxcpg_2025,mirsky2023vulchecker,fu2021linevul,debenedetti2024agentdojo}。

\subsection{结构化叙述策略}
本文采用结构化叙述综述而非严格 PRISMA 元分析。分析流程为：先按方法家族归类，再按审计对象和证据输入方式细分，最后按可复现性、可解释性、成本与部署约束比较。为降低叙述综述的主观性，本文对每类方法都提出同一组问题：依赖什么证据、解决什么审计子任务、在哪些样本上表现稳定、在哪些边界条件下失败、是否能接入可复现评测流水线。

\subsection{文献统计分析}
表~\ref{tab:literature_stats} 给出本文引用文献的分布统计。文献主要来自安全领域顶级会议（USENIX Security 52 篇、NDSS 39 篇、NeurIPS 23 篇）和软件工程/自然语言处理会议（ACL 9 篇、ICSE 5 篇、EMNLP 5 篇），以及 arXiv 预印本（26 篇）。时间分布上，2024--2026 年文献占比超过 85\%，反映该领域的快速发展。

\begin{table}[htbp]
\caption{引用文献分布统计}
\label{tab:literature_stats}
\centering
\small
\begin{tabularx}{\linewidth}{@{} p{0.35\linewidth}X @{}}
\rowcolor{tabhead}
\tabheader{来源} & \tabheader{数量} \\
\rowcolor{tabalt} USENIX Security & 52 \\
NDSS & 39 \\
\rowcolor{tabalt} NeurIPS & 23 \\
arXiv & 26 \\
\rowcolor{tabalt} ACL & 9 \\
IEEE S\&P & 7 \\
\rowcolor{tabalt} ICSE & 5 \\
\rowcolor{tabalt} EMNLP & 5 \\
ACM (TOSEM/CSUR) & 4 \\
\rowcolor{tabalt} 其他 (Springer, CEUR, etc.) & 18 \\
\hline
\rowcolor{tabhead} \tabheader{总计} & \tabheader{158} \\
\end{tabularx}
\end{table}

\subsection{证据追踪与可信度标注}
本文把关键结论绑定到四类证据：文献证据、本仓库 \path{benchmarks/audit-100/} 中的样本证据、\path{audit/audit-results/} 中的评测证据，以及 \path{docs/audit/audit-eval-iteration-log.md} 中的过程证据。表~\ref{tab:trust} 给出可信度标签。

\begin{table}[htbp]
\caption{结论可信度标签}
\label{tab:trust}
\centering
\small
\begin{tabularx}{\linewidth}{@{} p{0.10\linewidth}X X @{}}
\rowcolor{tabhead}
\tabheader{标签} & \tabheader{含义} & \tabheader{使用条件} \\
\rowcolor{tabalt} \textbf{高} & 文献、benchmark 样本和评测结果相互支持 & 至少代表性文献+样本对照+评测记录 \\
\textbf{中} & 有文献或样本支持，缺完整复现 & 可作为设计依据，不宜单独作强结论 \\
\rowcolor{tabalt} \textbf{低} & 主要来自方法推断或趋势观察 & 只能作为研究假设，需实证验证 \\
\end{tabularx}
\end{table}

\section{方法家族与代表性工作}

图~\ref{fig:taxonomy} 给出本文归纳的十类方法分类体系。方法按审计对象和技术路线分为四大类：代码级分析、结构化增强、系统级安全和评估治理。

\begin{figure}[htbp]
\centering
\begin{tikzpicture}[
  root/.style={draw=blue!60!black,fill=blue!15,rectangle,rounded corners=6pt,minimum width=4.5cm,minimum height=1cm,font=\normalsize\bfseries,text=blue!60!black,drop shadow},
  cat/.style={draw=#1!60!black,fill=#1!12,rectangle,rounded corners=4pt,minimum width=3cm,minimum height=0.9cm,font=\small\bfseries,text=#1!60!black,drop shadow},
  leaf/.style={draw=gray!50,fill=gray!5,rectangle,rounded corners=3pt,minimum width=2.6cm,minimum height=0.6cm,font=\footnotesize,text=black!80,drop shadow},
  edge/.style={->,>=latex,thick,color=gray!60,line width=0.8pt}
]

% Root
\node[root] (root) at (0,0) {LLM 辅助应用审计};

% Four categories
\node[cat=teal] (c1) at (-6,-2.5) {I. 代码级分析};
\node[cat=orange] (c2) at (-2,-2.5) {II. 结构化增强};
\node[cat=red] (c3) at (2,-2.5) {III. 系统级安全};
\node[cat=purple] (c4) at (6,-2.5) {IV. 评估治理};

% Edges from root
\draw[edge] (root.south) -- ++(0,-0.5) -| (c1.north);
\draw[edge] (root.south) -- ++(0,-0.5) -| (c2.north);
\draw[edge] (root.south) -- ++(0,-0.5) -| (c3.north);
\draw[edge] (root.south) -- ++(0,-0.5) -| (c4.north);

% Category I leaves
\node[leaf] (l1) at (-7.3,-4.2) {A. 直接代码审查};
\node[leaf] (l2) at (-6,-5.2) {B. 漏洞检测/修复};
\node[leaf] (l3) at (-4.7,-4.2) {C. 静态分析协同};
\draw[edge] (c1.south) -- ++(0,-0.3) -| (l1.north);
\draw[edge] (c1.south) -- ++(0,-0.3) -| (l2.north);
\draw[edge] (c1.south) -- ++(0,-0.3) -| (l3.north);

% Category II leaves
\node[leaf] (l4) at (-2.7,-4.2) {D. RAG/结构化};
\node[leaf] (l5) at (-1.3,-4.2) {F. Prompt/微调};
\draw[edge] (c2.south) -- ++(0,-0.3) -| (l4.north);
\draw[edge] (c2.south) -- ++(0,-0.3) -| (l5.north);

% Category III leaves (2 rows × 3)
\node[leaf] (l6) at (0.7,-4.0) {E1. IPI 攻防};
\node[leaf] (l7) at (2,-4.0) {E2. 越狱攻防};
\node[leaf] (l8) at (3.3,-4.0) {E3. Agent 安全};
\node[leaf] (l9) at (0.7,-5.0) {E4. 模型完整性};
\node[leaf] (l10) at (2,-5.0) {E5. LLM 安全测试};
\node[leaf] (l11) at (3.3,-5.0) {E6. 代码生成安全};
\draw[edge] (c3.south) -- ++(0,-0.2) -| (l6.north);
\draw[edge] (c3.south) -- ++(0,-0.2) -| (l7.north);
\draw[edge] (c3.south) -- ++(0,-0.2) -| (l8.north);
\draw[edge] (c3.south) -- ++(0,-0.5) -| (l9.north);
\draw[edge] (c3.south) -- ++(0,-0.5) -| (l10.north);
\draw[edge] (c3.south) -- ++(0,-0.5) -| (l11.north);

% Category IV leaves
\node[leaf] (l12) at (6,-4.2) {G. Benchmark/协议};
\draw[edge] (c4.south) -- (l12.north);

% Legend
\node[font=\tiny,text=gray,anchor=west] at (-7.5,-6.0) {图例：\textcolor{teal}{■}代码级 \textcolor{orange}{■}结构化 \textcolor{red}{■}系统安全 \textcolor{purple}{■}评估};

\end{tikzpicture}
\caption{LLM 辅助应用程序审计方法分类体系}
\label{fig:taxonomy}
\end{figure}

\subsection{代码审查的演进：从人工到 LLM 辅助}

\subsubsection{传统代码审查基础}
现代代码审查（Modern Code Review, MCR）的研究始于 Bacchelli 和 Bird~\cite{bacchelli_bird_2013} 在 ICSE 2013 的开创性工作：873 名开发者（44\% 响应率）和 165 名经理（28\% 响应率）共同表明，缺陷查找是首要动机，但实际产出更偏向知识传递与团队意识；在 570 条评论中，代码改进占 29\%，而 91\% 的开发者认为文件所有者参与审查有帮助、82\% 认为审查整体有价值。McIntosh 等~\cite{mcintosh_2015_impact} 进一步证明轻量级代码审查对软件质量的影响因项目而异：Qt 5.0/5.1、VTK、ITK 的缺陷率分别为 19\%/14\%/9\%/11\%，但即使有 100\% 覆盖，仍有大量后续缺陷暴露，说明覆盖率本身不足以保证质量。Kononenko 等~\cite{kononenko_2016_review_quality} 对 88 名 Mozilla 开发者的调查显示，审查质量的核心是代码质量与 thorough feedback：49\% 的编码引述指向代码质量、38\% 指向详尽反馈，且 100\% 的受访者认为对代码上下文的熟悉度重要。

在安全视角上，di Biase 等~\cite{dibiase_2016_security_chromium} 对 Chromium 项目中 132,000 条审查评论进行分析后发现，安全相关评论仅约占 1\%，在 9,765 条候选安全担忧评论中仅 71 条被人工确认为安全评论；C++ catastrophes 和 buffer overrun 是最常被遗漏的类别，而 command injection 从未在审查中被发现。Sadowski 等~\cite{sadowski_2018_google} 在 Google 对 9 million reviews 的案例研究展示了审查已服务于缺陷检测之外的多重目的：开发者每周中位数仅审查 4 个变更、初始反馈中位延迟不足 1 小时、总审查中位延迟不足 4 小时，且 70\% 的变更在 24 小时内合入。MacLeod 等~\cite{macleod_2018_trenches} 在 Microsoft 的 4,300 人调查与 18 次访谈中揭示，作者端最主要痛点是等待响应，审查端最主要痛点是时间不足。Davila 和 Nunes~\cite{davila_nunes_2021_slr_taxonomy} 对 139 篇论文的系统综述进一步确认：MCR 研究主要依赖开源项目、平均 1--2 名审查者/变更，且静态分析警告只有 6\%--22\% 会真正进入讨论，安全问题在审查中的可见性仍然偏低。

\subsubsection{LLM 时代的代码审查}
随着 LLM 的兴起，代码审查自动化进入了新阶段。早期尝试使用神经机器翻译~\cite{tufano_2021_towards} 和预训练模型~\cite{li_2022_codereviewer,tufano_2022_pretrained} 进行审查自动化，取得了可行性验证但精度有限。Yu 等~\cite{yu_2023_security_defect} 在 OpenStack 和 Qt 社区的研究发现，安全缺陷在代码审查中未被广泛识别，分歧和优先级降低是安全缺陷未被解决的首要原因。

LLM 时代的代码审查呈现三个特征：第一，多智能体协作架构的兴起。Tang 等~\cite{tang_2024_codeagent} 提出 CodeAgent，利用多个通信智能体协作检测提交信息不一致、漏洞引入和代码风格违规，在 3,545 个真实提交上进行了评估。Rasheed 等~\cite{rasheed_2024_ai_powered} 提出四专家智能体框架，分别关注代码异味、缺陷和改进建议。第二，工业部署的落地经验。Google 的 AutoCommenter~\cite{vijayvergiya_2024_autocommenter}（AIware 2024）已部署到数万开发者，覆盖 C++、Java、Python 和 Go，评论解决率 40\%、60\% 评论被认为有用，采用 per-URL 阈值和 beam search 覆盖 68\% 最佳实践 URL。Google 的 DIDACT V2 系统~\cite{frommgen_2024_resolving}（ICSE-SEIP 2024）通过 ML 辅助解决审查评论，每年自动解决 7.5\% 的审查评论（Java 9.5\%、C++ 7.5\%），节省数十万工程师小时。Microsoft 的 WhoDo~\cite{asthana_2024_whodo} 在生产环境中部署了审查者推荐系统。第三，LLM 审查的不完美性。Olewicki 等~\cite{olewicki_2024_llm_review_impact} 在 Mozilla+Ubisoft 的首个活跃用户研究（59 个审查者、587 个补丁）中发现接受率仅 8.1\%/7.2\%，但 23\%/28.3\% 的评论被认为有价值——接受率低不等于无用。Cihan 等~\cite{cihan_2025_llm_review_eval} 发现 GPT-4o 和 Gemini 2.0 Flash 在 492 个 AI 生成代码块上的正确分类率分别仅为 68.5\% 和 63.9\%。Pornprasit 和 Tantithamthavorn~\cite{pornprasit_2024_gpt35} 发现 few-shot learning 可将 GPT-3.5 性能提升至少 46\%，但 persona 提示反而降低性能。

Jin 和 Chen~\cite{overcorrection_2026} 展示了 LLM 代码审查器在需求一致性判断中容易出现"纠错过度"的偏差。Tufano 等~\cite{tufano_2024_strengths_weaknesses} 在 2,291 个预测上使用 120 标签分类学评估，发现 T5 代码审查准确率仅 14.08\%/2.11\%，且\textbf{训练数据集中约 25\% 为噪声}——揭示了代码审查自动化的核心矛盾：ChatGPT 难以像人类审查者那样评论代码，专用 DL 模型在审查任务上优于通用 LLM，但数据质量是根本瓶颈。Liu 等~\cite{liu_2025_noisy_data} 和 Auto Code Review~\cite{auto_code_review_2024} 进一步证明开源代码审查数据集中包含大量噪声评论，即使经过启发式和 ML 清洗方法仍持续存在，导致模型生成低质量输出。Hong 和 Baik~\cite{hong_2025_rag_review} 通过 RAG 结合生成式和检索式方法，缓解生成模型倾向产生高频通用 token 的问题。

在安全审查方面，Peng 等~\cite{peng_2025_icodereviewer} 的 iCodeReviewer 使用混合提示方法实现 63.98\% F1 的安全问题识别率和高达 84\% 的审查评论接受率。Charoenwet 等~\cite{charoenwet_2024_sast_secure_review} 的实证研究发现，单个 SAST 工具可在 52\% 的漏洞代码变更中标记警告，但 76\% 的漏洞函数中的警告不相关，22\% 的漏洞贡献提交仍未被检测。Watanabe 等~\cite{watanabe_2024_chatgpt} 的 229 条 ChatGPT 审查评论研究发现，审查者将 ChatGPT 用于外包实现、重构、修复和测试任务，但对审查深度和准确性仍有担忧。Yang 等~\cite{yang_2026_roadmap} 提出了现代代码审查的十年路线图， envisioning 三个范式转变：从被动助手到 IDE 原生主动 AI 协作者、从准确率指标到价值驱动评估、从自主工具到互惠人机反馈循环。

\subsection{漏洞检测与修复}

\subsubsection{深度学习时代的漏洞检测}
漏洞检测经历了从规则引擎到深度学习再到 LLM 的演进。Kim 等~\cite{kim2017vuddy} 的 VUDDY 提出了可扩展的漏洞代码克隆检测方法，结合函数级抽象和代码相似度在大规模 OSS 生态中发现传播漏洞，实现近零误报。Li 等~\cite{li2018vuldeepecker} 的 VulDeePecker（NDSS 2018）是首个基于深度学习的漏洞检测系统，在 61,638 个代码片段上达到 TPR=93.0\%、FPR=5.7\%、F1=90.5\%（远超 Flawfinder 27.7\% 和 RATS 20.2\%），并发现 4 个此前未知的漏洞。Zhou 等~\cite{zhou2019devign} 的 Devign（NeurIPS 2019）在代码属性图上构建图神经网络，比基线方法平均准确率提升 +10.51\%、F1 提升 +8.68\%，对最新 40 个 CVE 的准确率达 74.11\%。Feng 等~\cite{feng2020codebert} 的 CodeBERT 提供了通用代码表示的预训练基线。Fu 和 Tantithamthavorn~\cite{fu2021linevul} 的 LineVul（MSR 2022）首次将 Transformer 应用于行级漏洞预测，在 188,000+ 函数上达到函数级 F-measure=0.91（比基线高 160--379\%），行级 Top-10 准确率 0.65，对 Top-25 最危险 CWE 类型准确率 75--100\%。Ahmad 等~\cite{ahmad2021plbart} 的 PLBART 通过去噪自编码在大规模代码上预训练，为漏洞解释和修复生成提供了基础。

近期工作进一步深化了漏洞检测的技术路线。Mirsky 等~\cite{mirsky2023vulchecker} 的 VulChecker（USENIX Security 2023）在 19 个真实项目中检测到 24 个 CVE（商业工具 Helix QAC 仅检测到 4 个，精度提升 6$\times$），并发现一个可利用的零日漏洞。Peng 等~\cite{peng2023ptlvd} 的 PTLVD 将程序切片与 Transformer 结合实现行级检测。Ding 等~\cite{ding2023vulhawk} 的 VulHawk 实现了跨架构（x86、ARM、MIPS）二进制代码漏洞搜索。Woo 等~\cite{woo2023v1scan} 的 V1SCAN（USENIX Security 2023）结合版本匹配和代码相似度，比现有方法多发现 50\% 的漏洞（137 个），精度 96\%、召回 91\%，将 FPR 从 71\% 降至 4\%。Shao 等~\cite{shao2024fvddpm} 的 FVD-DPM（USENIX Security 2024）将漏洞检测重构为条件扩散概率生成过程，语句级定位 IoU 比 Cppcheck 高 +37\%，比 DeepLineDP 高 +42\%。Shimmi 等~\cite{shimmi2024vulsim} 利用多维邻居嵌入改善跨漏洞类型检测。Liu 等~\cite{liu2024contrastive} 使用对比学习解决漏洞 CWE 分类的长尾分布问题。Wang 等~\cite{wang2024taskfusion} 提出面向任务的知识融合，针对不同代码区域选择性整合异构表示。Chen 和 Xie~\cite{chen2024vullibgen} 利用 LLM 生成受漏洞影响包名称。

\subsubsection{LLM 漏洞检测的能力边界}
多项实证研究揭示了 LLM 在漏洞检测中的系统性局限。Paquette 等~\cite{paquette2024llmscannot} 在 IEEE S\&P 2024 引入 SecLLMHolmes 框架，证明 LLM 无法可靠区分漏洞与非漏洞代码，也不能一致地推理安全属性。Risse 和 Böhme~\cite{risse2024limits} 在 USENIX Security 2024 提出 VulnPatchPairs 基准（26,200 个 C 函数及其补丁对），发现六个顶级 ML4VD 模型在区分漏洞函数与其修复补丁时准确率\textbf{低于随机猜测}（$<50\%$），且数据增强仅使模型过拟合到特定变换——当训练集和测试集使用不同的语义保持变换时，准确率恢复仅 2.5--5.9\%。Liu 等~\cite{liu2024eatvul} 展示了 ChatGPT 可生成语义保持的对抗代码变体来逃避 DL 漏洞检测器。

Niu 等~\cite{niu2025largetomammoth} 在 NDSS 2025 对 5 个 LLM 家族、21 个模型（7B--70B 参数）的大规模评估揭示：(1) Gemma 7B 在 Java 零样本上达到最高 F1（57.98\%），而 LLaMA-3 70B 的 F1 仅 5.10\%，\textbf{新一代模型反而不如旧版}；(2) GPT-4 的准确率（37.86\%）低于开源 Gemma（78.57\%）；(3) 上下文窗口扩大一致提升性能（2048$\to$8192 tokens 使 LLaMA-3 8B 准确率从 21.62\% 升至 52.70\%），但 few-shot learning \textbf{反而降低}性能（Gemma 7B F1 从 57.98\% 降至 44.60\%）；(4) 量化效果因模型而异，无统一最优策略。Khare 等~\cite{khare2025llmvulndet} 发现 LLM 在已知漏洞模式上表现尚可，但在复杂、多步或应用特定漏洞上有效性显著下降。Rossi 等~\cite{rossi2025crossdomain} 发现跨领域迁移时性能显著退化。Zhang 等~\cite{zhang2024prompt} 证明精心设计的 prompt 可显著改善 ChatGPT 零样本漏洞检测——链式思维将 C/C++ 准确率从 52.3\% 提升至 74.1\%，辅助信息（数据流图）将 Java 准确率提升至 74.7\%。Song 等~\cite{song2025nodemedic} 开发了 Node.js 包的细粒度动态分析。

\subsubsection{自动化漏洞修复}
漏洞修复是检测的自然延伸。Zhou 等~\cite{zhou2025vulndetrepair} 的 ACM TOSEM 综述回顾 58 项研究、37 个 LLM，发现最优检测准确率仅 67.6\%、修复准确率仅 $\sim$20\%，且 63\% 修复方法依赖微调、37\% 依赖 prompt，6 项关键限制包括函数级输入粒度、缺乏过程间分析和 IDE 集成不足。Zhang 等~\cite{zhang2025tplpatch} 提出第三方库 1-day 漏洞自动补丁方法。USENIX Security 2025 出现三项重要修复工作：Wen 等~\cite{wen2025appatch} 的 AppAtch 使用自适应提示迭代优化 LLM 修复，在零日漏洞上 F1 达 36.46\%，无需利用证据即可超越传统 APR 工具；Kim 等~\cite{kim2025san2patch} 的 SAN2PATCH 仅用 sanitizer 日志和源码结合思维树推理，在真实漏洞上达到 79.5\% 正确补丁率；Yu 等~\cite{yu2025patchagent} 的 PatchAgent 实现端到端自主修复，修复成功率 92.13\%，从定位到验证无需人工干预。Zhou 等~\cite{zhou_vuln_repair_2024} 强调"发现漏洞"和"生成补丁"不是同一件事。SoK~\cite{sok2025avr} 系统化整理了自动修复全景，将方法分为规则、学习和 LLM 三类，识别了过度依赖合成 benchmark 的评估陷阱。

\subsubsection{综合综述}
近期多篇综述系统化了 LLM 安全检测领域。Kaniewski 等~\cite{kaniewski2026detecting} 的 ACM TOSEM 综述分析了 263 篇研究（2020--2025），是该领域最大规模的 SLR，发现研究高度碎片化——各研究使用自定义数据集、指标和数据划分，难以跨研究比较。Sheng 等~\cite{sheng2026llmsecurity} 的 ACM CSUR 综述涵盖 80+ 篇论文和 33 个不同 LLM，发现 GPT-4 是最常用模型（29 次使用），系统化了提示策略、微调方法和评估方法论，识别了幻觉、数据污染和可复现性等关键挑战。Chen 等~\cite{chen2026lmcode} 的 ACM TOSEM 综述系统回顾了 68 篇关于 CodeLM 安全性的论文，涵盖后门攻击和对抗攻击两大类别。Zhang 等~\cite{sheng2025cybersecurity_slr} 的综述涵盖 300+ 项工作、25 个 LLM 和 10+ 下游场景。各综述一致指出：高质量数据集稀缺是首要瓶颈，混合方法（LLM + 传统分析）最具前景，仓库级分析仍未解决。

\subsection{LLM 与静态分析协同}
静态分析协同是当前最具工程可行性的路线之一。其基本思想是：静态规则、污点分析或图分析负责提供可重复、可解释的候选证据，LLM 负责判断这些证据在当前上下文中是否真正构成安全威胁。相关研究包括 LLM 辅助程序分析综述、静态分析与 LLM 协同、CPG 引导的漏洞检测，以及 LLM 与静态工具的系统 benchmark~\cite{program_analysis_survey_2025,synergizing_static_llm_2025,llmxcpg_2025,llm_vs_static_2025}。

该路线适合应用程序审计，因为审计同时需要高召回和低误报。静态工具负责发现可疑点，LLM 负责解释为什么可疑或为什么不危险。优势是可追溯、可审计、可接入 CI；局限是前处理链条更长，且静态证据缺失会使 LLM 只能在残缺上下文中判断。Charoenwet 等~\cite{charoenwet_2024_sast_secure_review} 的实证数据支持了这一判断：SAST 工具对漏洞代码变更的覆盖率仅 52\%，且 76\% 的警告不相关，说明 LLM 在过滤误报方面有显著价值。

例如，LLMxCPG~\cite{llmxcpg_2025} 以 CPG 为中心组织程序上下文，利用图结构和切片把漏洞候选压缩成更高信噪比的证据；Synergizing~\cite{synergizing_static_llm_2025} 则直接把"静态分析 + LLM"组合成混合工作流；LLM vs Static~\cite{llm_vs_static_2025} 在 10 个 C\# 项目（63 个嵌入漏洞）上的系统比较证实二者互补而非替代：GPT-4.1 F1=0.797、Mistral Large F1=0.753、DeepSeek V3 F1=0.750，远超 SonarQube（F1=0.260）、CodeQL（F1=0.386）和 SnykCode（F1=0.546）——LLM 优势来自更宽的代码上下文召回，但所有 LLM 因 tokenization 在行/列定位上存在偏差；推荐混合流水线：LLM 负责早期广域分诊，确定性扫描器负责高保证验证。Kaniewski 等~\cite{kaniewski2026detecting} 在工业流水线评估中发现，LLM 在定位和解释方面表现优秀，但严重性校准仍需人工确认。

\subsection{检索增强与结构化表示增强}
结构化/RAG 路线的核心思想是：不要把整份源码直接交给模型，而是先将程序知识压缩为适合审计的证据，再让 LLM 在证据上推理。常见输入包括 CPG、CFG、DFG、调用链、数据流路径、敏感源/汇点、API 文档、历史提交、相似漏洞案例和项目规范~\cite{program_analysis_survey_2025,synergizing_static_llm_2025,llmxcpg_2025}。

结构化表示的收益是缓解上下文窗口限制、提升证据密度、让模型围绕"为什么危险"而不是"代码长什么样"判断；局限是图构建和切片本身可能出错，关键边或节点缺失会直接影响结论。RAG 的收益是补充 API 语义和项目约定；风险是检索到的外部文本可能含有噪声、误导说明或提示注入。Hong 和 Baik~\cite{hong_2025_rag_review} 通过 RAG 缓解代码审查中生成模型的通用 token 倾向。Tan 等~\cite{tan2026prarag} 的 PRA-RAG 提供了首个可证明鲁棒的 RAG 聚合机制，通过约束被污染检索文档对最终输出的影响，将攻击成功率从基线的 >40\% 降至 $\sim$1\%，同时在干净数据上保持 $\sim$71\% 准确率——是目前唯一同时提供理论保证和实验验证的 RAG 防御方案。

\begin{table}[htbp]
\caption{结构化/RAG 审计可靠性评估维度}
\label{tab:rag_reliability}
\centering
\small
\begin{tabularx}{\linewidth}{@{} p{0.14\linewidth}X X @{}}
\rowcolor{tabhead}
\tabheader{层次} & \tabheader{典型失败模式} & \tabheader{对结论的影响} \\
\rowcolor{tabalt} 表示构建 & 动态导入、反射导致调用边缺失 & 恶意链被切断，形成漏报 \\
证据切片 & 只保留危险 API，不保留校验条件 & 误报或无法证明可利用性 \\
\rowcolor{tabalt} 检索质量 & 相似但无关案例进入上下文 & LLM 被错误先验牵引 \\
检索投毒 & 对抗性文档注入恶意指令 & IPI，审计结论被篡改 \\
\rowcolor{tabalt} 证据冲突 & 文档声称安全，代码路径显示外传 & 应升级为 UNCERTAIN 或人工复核 \\
可复现性 & 检索源、摘要、排序变化 & 指标不可比较，结论难复核 \\
\end{tabularx}
\end{table}

\subsection{面向智能体与应用执行链的审计}
随着 LLM 从代码生成器扩展为 agent，审计对象不再只是静态源码，而是会检索、调用工具、规划行动并在多轮交互中改变状态的执行系统。2024--2026 年间，agent 安全已成为独立研究领域，涵盖 prompt injection 攻防、jailbreak 攻防、agent 行为链安全和模型完整性等多个方向。图~\ref{fig:agent_security} 给出 agent 安全的攻防全景。

\begin{figure}[htbp]
\centering
\begin{tikzpicture}[
  attack/.style={draw=red!50!black,fill=red!10,rectangle,rounded corners=4pt,minimum width=3cm,minimum height=0.8cm,font=\small,align=center,text=red!60!black,drop shadow},
  defense/.style={draw=green!50!black,fill=green!10,rectangle,rounded corners=4pt,minimum width=3cm,minimum height=0.8cm,font=\small,align=center,text=green!50!black,drop shadow},
  target/.style={draw=blue!50!black,fill=blue!10,circle,minimum size=2.5cm,font=\small\bfseries,text=blue!60!black,align=center,drop shadow},
  attackarrow/.style={->,>=latex,line width=1.5pt,color=red!50!black},
  defensearrow/.style={->,>=latex,line width=1.5pt,color=green!50!black}
]

% Central agent
\node[target] (agent) at (0,0) {LLM Agent\\{\scriptsize 审计对象}};

% Attack vectors (left)
\node[attack] (a1) at (-5.5,3) {\textbf{IPI} 间接提示注入};
\node[attack] (a2) at (-5.5,1.5) {\textbf{越狱} Jailbreak};
\node[attack] (a3) at (-5.5,0) {\textbf{恶意技能} Skills};
\node[attack] (a4) at (-5.5,-1.5) {\textbf{后门} Backdoor};
\node[attack] (a5) at (-5.5,-3) {\textbf{泄露} Prompt/KV};

% Defenses (right)
\node[defense] (d1) at (5.5,3) {Spotlighting, StruQ};
\node[defense] (d2) at (5.5,1.5) {JailbreakBench, RPO};
\node[defense] (d3) at (5.5,0) {IsolateGPT, TaskShield};
\node[defense] (d4) at (5.5,-1.5) {LMSanitator, Barbie};
\node[defense] (d5) at (5.5,-3) {执行隔离, 审计日志};

% Attack arrows
\foreach \a in {a1,a2,a3,a4,a5} {
  \draw[attackarrow] (\a.east) -- (agent.west);
}

% Defense arrows
\foreach \d in {d1,d2,d3,d4,d5} {
  \draw[defensearrow] (\d.west) -- (agent.east);
}

% Column headers
\node[font=\small\bfseries,text=red!50!black,anchor=south] at (-5.5,4) {⚔ 攻击向量};
\node[font=\small\bfseries,text=green!50!black,anchor=south] at (5.5,4) {🛡 防御机制};

% Key findings box
\node[font=\small,draw=orange!40,fill=orange!5,rounded corners=4pt,text width=12cm,align=center,text=black!70] at (0,-4.5)
  {\textbf{关键发现}：能力越强越脆弱（Claude PNA=90.75\% 但 IPI ASR=94.5\%）~\cite{zhang2024asb}；越狱统计意义上不可避免~\cite{su2024missionimpossible}；防御降 15--20\% 效用~\cite{debenedetti2024agentdojo}；微调模型比 prompt 模型鲁棒 6$\times$~\cite{zhan2024injecagent}};

\end{tikzpicture}
\caption{Agent 安全攻防全景图}
\label{fig:agent_security}
\end{figure}

\subsubsection{间接提示注入攻防}
间接提示注入（IPI）是 agent 安全的核心威胁。Zhan 等~\cite{zhan2024injecagent} 的 InjecAgent（ACL 2024）引入 1,054 个测试用例的 benchmark（17 用户工具 $\times$ 62 攻击工具），发现 ReAct-prompted GPT-4 ASR 为 24\%（增强攻击下升至 47\%），Llama2-70B 超过 80\%，但\textbf{微调模型显著更鲁棒}——微调 GPT-4 ASR 仅 3.8\%，远低于 prompt 版本的 24\%。Liu 等~\cite{liu2024promptinjbench}（USENIX Security 2024）提出统一形式化框架，在 10 个 LLM、7 个任务上评估 5 种攻击和 10 种防御——其 Combined Attack 在 10 个 LLM 上平均 ASV=0.62、MR=0.78，当目标任务和注入任务类型相同时检测极其困难，没有单一防御普遍有效。Li 等~\cite{li2024instrfollow} 揭示了一个根本性张力：指令跟随能力越强的模型反而越容易被提示注入，因为它们更倾向于遵从注入指令而非原始指令。Kong~\cite{kong2024injectbench} 开发了系统化的 IPI benchmark 框架。Chen 等~\cite{chen2025canipi} 发现检测分类器在独立设置中准确率高，但在端到端 agent 场景中因分布偏移而显著退化。

防御方面，Hines 等~\cite{hines2024spotlighting} 的 Spotlighting 通过不可见标记区分可信指令与不可信数据——datamarking 将 GPT-3.5 摘要任务 ASR 从 50\% 降至 3.10\%，编码方案降至 0\%，且不损害 NLP 任务性能。Ayub 和 Majumdar~\cite{ayub2024embedding} 在 467,057 条 prompt 上证明基于嵌入的轻量分类器可以高准确率检测提示注入——Random Forest + OpenAI 嵌入达 AUC=0.764、F1=0.868，在 AUC 和精度上超越 4 个 DeBERTa SOTA 分类器（DeBERTa 召回 >0.91 但精度仅 0.75--0.77）。Chen 等~\cite{chen2025struq} 的 StruQ 在 API 层面强制 prompt 与数据分离，对 Naive/Ignore/Completion-Real 等攻击实现 ASR 从 96\% 降至 0\%，TAP 从 97\% 降至 9\%，同时保持 AlpacaEval 效用不变（67.2\%→67.6\%）。Jia 等~\cite{jia2025taskshield} 的 Task Shield 通过任务对齐约束阻止 IPI 导致的偏离。Wen 等~\cite{wen2025instrdetect} 的 InstructDetector 利用 LLM 行为状态（隐藏层状态+梯度）检测 IPI 攻击，仅需 200 个训练样本即达 99.60\% 域内精度和 96.90\% 域外精度，在 BIPIA benchmark 上将 GPT-3.5 ASR 从 33.57\% 降至 0.12\%、GPT-4o 从 39.68\% 降至 0.13\%。Chen 等~\cite{chen2025defenseleverage} 用攻击技术反制攻击，通过对抗性增强训练提高鲁棒性。Tan 等~\cite{tan2026prarag} 的 PRA-RAG 提供了首个可证明鲁棒的 RAG 聚合机制。Ramakrishnan 和 Balaji~\cite{ramakrishnan2025pidefense} 提出多层防御框架，将 IPI 成功率降低超 70\%。

Evertz 等~\cite{evertz2024whispers} 揭示 LLM 集成系统中的机密性威胁：Llama 3.1 在无防护时 99\% 泄露密钥，显式"不要泄露"指令仍泄露 21--34\%，越狱恢复至 94--97\%；Context Ignoring 是最有效攻击（8b 上 91\% ASR）；工具集成可将提取成功率提升 73\%。Kim 等~\cite{kim2025manyshot} 证明上下文长度（而非 shot 数）是 many-shot 攻击有效性的主要决定因素。Evtimov 等~\cite{evtimov2025wasp} 的 WASP benchmark（84 个对抗任务）展示了网页导航中 prompt injection 可劫持商业 agent——中间劫持成功率高达 85.7\%（o1 模型），但端到端目标完成率仅 0--16.7\%，揭示"安全源于无能"现象：攻击可劫持 agent 但难以完成复杂攻击目标。URL 注入比纯文本注入有效 2.6$\times$（GPT-4o 上 61.9\% vs 23.8\%），防御性系统提示可将 ASR 从 42.9\% 降至 22.6\%。

\subsubsection{越狱攻击与防御}
越狱攻击旨在绕过 LLM 的安全对齐。Deng 等~\cite{deng2024masterkey} 的 MasterKey（NDSS 2024）利用基于时间的防御探测技术逆向工程聊天机器人防御机制，自动生成越狱 prompt，平均成功率 21.58\%（基线的 3$\times$），并发现越狱 prompt 具有 GPT 专属性——对 Bard 和 Bing Chat 的迁移率接近零。Yu 等~\cite{yu2024dontlisten} 对野外越狱 prompt 进行大规模实证研究，分类出角色扮演、对抗后缀和逻辑推理链三种核心攻击机制。Yu 等~\cite{yu2024fuzzer} 将软件测试的 fuzzing 方法引入越狱发现，识别出数千个未知攻击向量。

NeurIPS 2024 集中发表了多项越狱研究。Xu 等~\cite{xu2024bagoftricks} 将越狱分解为 token 级和 prompt 级组件，发现语义混淆和角色扮演框架主导字符级扰动。Hu 等~\cite{hu2024adc} 的 ADC 实现从密集到稀疏的 token 级优化，查询效率远超先前方法。Zheng 等~\cite{zheng2024fewshot} 证明精心构造的 few-shot 演示可跨模型迁移并绕过已知防御。Su 等~\cite{su2024missionimpossible} 从统计学习角度证明越狱在理论上不可避免：任何有限训练集上的对齐都留有可利用的残余漏洞。

USENIX Security 2025 继续推进：Gong 等~\cite{gong2025papillon} 的 PAPILLON 结合 fuzz testing 和 LLM 引导变异生成语义连贯的越狱 prompt，对 GPT-3.5 实现 >90\% ASR、GPT-4 >80\%、Gemini-Pro >74\%，生成 prompt 的困惑度仅 34.61（远低于检测阈值 58.83），绕过 perplexity-based 检测器的能力显著优于基线。Krauss 等~\cite{krauss2025twinbreak} 的 TwinBreak 同时嵌入良性和恶意请求，利用 LLM 处理第二条指令的倾向，在 16 个模型（1B--72B 参数）上实现 89--98\% ASR，运行时间 <5 分钟，工具损失 $\leq$4.5\%。Zhang 等~\cite{zhang2025tasklevel} 识别了代码生成、翻译等任务级漏洞。

防御方面，Chao 等~\cite{chao2024jailbreakbench} 的 JailbreakBench 建立了标准化可复现评测协议，揭示许多已发表防御在标准化测试下提供的鲁棒性微不足道——其 Prompt+RS 基线对 4 个 LLM 分别实现 89/90/93/78\% ASR，其中 Erase-and-Check 是最强防御但仍难以完全阻断。Han 等~\cite{han2024wildguard} 的 WildGuard 发布开放审核工具包，联合执行恶意意图检测、响应安全分类和拒绝检测。Jiang 等~\cite{jiang2024wildteaming} 从真实交互中挖掘 5,700 个新越狱策略，对抗性训练产出更鲁棒模型而不降低有用性。Zhou 等~\cite{zhou2024rpo} 的 RPO 自动生成鲁棒指令前缀，提供 Lipschitz 有界威胁模型下的可证明鲁棒性保证。Gong 等~\cite{gong2025safetyMisalign} 揭示对齐技术（RLHF、DPO）产生"对齐税"——HarmBench ASR 从 42\% 降至 36.95\%、AdvBench 从 28.40\% 降至 20.89\%，但在某些维度反而降低安全性。

\subsubsection{Agent 系统安全}
Agent 安全审计关注"输入--决策--工具--输出"的完整链路。Debenedetti 等~\cite{debenedetti2024agentdojo} 的 AgentDojo（NeurIPS 2024）引入动态评估环境（97 个任务、629 个安全测试、4 个环境、74 个工具），agent 在不可信数据上执行真实工具调用——最佳 agent 的良性效用天花板仅 66\%，Slack 套件 ASR 高达 92\%，工具过滤防御可将 ASR 降至 7.5\% 但 17\% 测试用例仍然失败，所有防御均导致 15--20\% 效用损失。Zhang 等~\cite{zhang2024asb} 的 Agent Security Bench（ICLR 2025）覆盖 10 个场景、400+ 工具、13 个 LLM 骨干、13 种攻击和 11 种防御，发现平均 ASR 高达 84.30\%（PoT 后门攻击对 GPT-4o 达 100\%），防御平均有效性仅 46.91\%，且\textbf{能力越强的模型越脆弱}——Claude 3.5 Sonnet 效用最高（PNA=90.75\%）但 IPI ASR 也达 94.50\%。Zhang 等~\cite{zhang2026agentaudit} 的 Agent Audit 检查模型权重之外的工具代码、部署配置和权限设置。

Wu 等~\cite{wu2025isolateGpt} 的 IsolateGPT 提出 hub-and-spoke 执行隔离架构，每个应用运行在独立 spoke 中配备专用 LLM 实例，在 1,598 个攻击场景上将 VanillaGPT 的 81.3\% 跨应用攻击成功率大幅降低，75.73\% 查询性能开销 <30\%，同时保持 100\% 功能完整性。Ayzenshteyn 等~\cite{ayzenshteyn2025cloak} 的 CHeaT 框架使用 Stackelberg 博弈建模的欺骗（蜜罐）和反击策略，在 11 台 CTF 机器上实现 100\% 防御成功率（0/3 被攻破），累积 5 个随机陷阱后 DSR >95\%，利用 LLM 的训练偏差、DFS 搜索倾向和幻觉弱点。Liu 等~\cite{liu2025taintstyle} 的 AgentFuzz 将程序安全的污点分析和有向灰盒模糊测试适配到 LLM agent，在 20 个开源 agent 应用上发现 34 个 0-day 漏洞（23 个 CVE），精度 100\%、召回 97.14\%，比 LLMSmith 精度提升 33$\times$。Liu 等~\cite{liu2026malskills,malicious_skills_2026}（USENIX Security 2026）对 98,380 个社区 agent 技能的分析发现 157 个确认恶意技能（0.16\%），包含 632 个漏洞覆盖 6 个 MITRE ATT\&CK 杀伤链阶段，每个恶意技能平均 4.03 个漏洞，>50\% 源自单一威胁行为者。Agent Audit~\cite{agent_audit_2026} 提出 agent 感知静态分析框架，在 42 个标注漏洞上达到 F1=0.909（Semgrep 0.385、Bandit 0.458），召回率 95.24\%，比 Semgrep 高 4$\times$，MCP 漏洞覆盖率 100\%。Wu 等~\cite{wu2025promptleakage}（NDSS 2025）的 PromptPeek 利用 KV-cache 共享侧信道攻击，在多租户 LLM 服务中以 96--100\% 成功率提取其他用户的 prompt，GQA 架构使攻击更容易（100\% 全场景）。

Xu 等~\cite{xu2026duality,agent_security_duality_2026}（AI Review 2026）通过"双重性"框架综述 agent 安全（89 页），揭示仅 0.1\% 数据投毒即可实现 >80\% 后门攻击成功率且对良性性能影响 <1\%，即使最安全的 agent 仍有 23.9\% 测试用例存在失败风险。Khonde 等~\cite{khonde2026ragsecurity,rag_agent_security_2026} 提供 RAG agent 管道的端到端安全综述，涵盖检索投毒、IPI 和工具攻击三类威胁。Chen 等~\cite{chen2025canipi}（ACL 2025）构建 Inj-SQuAD/Inj-TriviaQA 基准，发现通用 LLM 检测 IPI 准确率仅 42--79\%，但专门训练的 DeBERTa 达 99.12\%，Segment-DeBERTa 组合过滤可 100\% 移除注入指令，将 StruQ 的 ASR 进一步降至 0.11\%。Ramakrishnan 和 Balaji~\cite{ramakrishnan2025pidefense} 在 847 个对抗测试上验证三层防御框架，将整体 ASR 从 73.2\% 降至 8.7\%，延迟开销仅 2.1\%。Prompt injection agents~\cite{prompt_injection_agents_2025} 强调间接提示注入可绕过主模型的表面安全。

\subsubsection{模型完整性与后门威胁}
当审计系统依赖微调模型或第三方预训练模型时，模型本身的完整性成为审计可信性的前提。后门攻击可在模型中植入触发器，使特定输入产生攻击者期望的输出，从而破坏审计判断。

2023--2025 年间，后门攻防研究呈现多样化演进。Beatrix~\cite{li2023beatrix} 使用 Gram 矩阵进行鲁棒后门检测。ASSET~\cite{pan2023asset} 在多重 DL 范式下实现鲁棒后门数据检测。PELICAN~\cite{zhang2023pelican} 揭示自然训练模型中可利用的后门在二进制代码分析中的存在。VILLAIN~\cite{bai2023villain} 针对垂直联邦学习发起后门攻击。Trojan Source~\cite{boucher2023trojan} 展示了不可见字符可产生源码级别的隐形漏洞。

与 LLM 审计直接相关的是：Wei 等~\cite{wei2024lmsanitator} 的 LMSanitator（NDSS 2024）在 960 个模型上达到 92.8\% 后门检测准确率，将 ASR 降至 <1\%，且仅降低 0.075\% 准确率（远优于 ONION 的 2.2\%），揭示 prompt-tuning 容易受到预训练模型中任务无关后门的影响。Zhang 等~\cite{zhang2024instrbackdoor}（USENIX Security 2024）展示针对定制化 LLM（如 OpenAI GPTs）的指令后门攻击，词级触发器在 6 个 LLM 上均达到 ASR=1.000，但防御端 5.8\% 误报率意味着 300 万个 GPTs 中约 17.4 万个会被误标。BackdoorAlign~\cite{backdooralign2024}（NeurIPS 2024）仅用 11 条安全样本即将 Llama-2 ASR 从 94.91\% 降至 3.64\%、GPT-3.5 从 75.64\% 降至 14.91\%，比基线防御节省 27$\times$ 数据量。Philosopher's Stone~\cite{philosophers2025} 展示 LLM 插件中的木马化风险。Yan 等~\cite{yan2024llmbackdoor} 进一步证明 LLM 可辅助对代码补全模型发动后门攻击，扩大后门注入的攻击面。

防御方面，Beatrix~\cite{li2023beatrix}（NDSS 2023）使用 Gram 矩阵在动态后门检测上达 F1=91.1\%，远超 SOTA 的 36.9\%。Barbie~\cite{barbie2025}（NDSS 2025）基于潜在可分性在 >10,000 个模型上验证，对 14 类后门攻击的 TPR 接近 100\%、FPR <3.65\%，比 FreeEagle 在样本特定攻击上提升 +43\% TPR。CLIBe~\cite{clibe2025} 检测 Transformer NLP 模型中的动态后门。DShield~\cite{dshield2025}（NDSS 2025）通过差异学习同时防御脏标签和干净标签 GNN 后门攻击，将 Cora 上 ASR 从 97.81\% 降至 0.74\%。这些工作共同表明：若审计管线中使用的 LLM 或嵌入模型本身被后门化，则所有下游审计结论均不可信。因此，高可靠审计系统必须将模型完整性验证纳入信任链。

\subsubsection{LLM 增强安全测试}
LLM 不仅被审计，也在主动参与安全测试。Deng 等~\cite{deng2024pentestgpt} 的 PentestGPT（USENIX Security 2024）将 LLM 用于自动化渗透测试，在 13 个目标上将任务完成率提升 228.6\%（GPT-4 完成 5/13 目标），picoMini CTF 排名 24/248，但上下文丢失是主要失败原因（195 次试验中 74 次）。Asmrita 等~\cite{asmrita2024busybox}（USENIX Security 2024）在 293 个 BusyBox ELF 二进制上利用 LLM 生成种子和崩溃重用技术，无需传统模糊测试即在最新 BusyBox 中发现 97 个崩溃（19 个唯一）。DeepGO~\cite{deepgo2024}（NDSS 2024）使用预测引导的有向灰盒模糊测试，到达目标站点速度比 AFLGo 快 3.23$\times$，暴露 19/20 个已知漏洞。EnclaveFuzz~\cite{enclavefuzz2024}（NDSS 2024）在 14/20 个 SGX enclave 中发现 162 个 bug，代码覆盖率 16.53\%（SGXFuzz 仅 4.57\%）。LLM 引导的协议模糊测试~\cite{llm_protocol_fuzz_2024} 从 RFC 规范自动生成结构化种子和状态感知变异策略。

在传统安全测试方面，Fuzztruction~\cite{bars2023fuzztruction}（USENIX Security 2023）通过对生成器而非输入进行故障注入，在复杂格式消费者上比 AFL++ 等工具平均多 21\% 覆盖率（部分目标达 70\%），发现 151 个唯一崩溃和 4 个 CVE。NAUTILUS~\cite{deng2023nautilus}（USENIX Security 2023）针对 RESTful API 的多步漏洞检测，比现有扫描器多发现 141\% 漏洞，在真实服务中发现 23 个零日漏洞（12 个 CVE），包括 Atlassian Confluence 远程代码执行和 Microsoft Azure 高危漏洞，揭示 82\% 的 API 漏洞需跨多个操作才能触发。ICSPatch~\cite{rajput2023icspatch}（USENIX Security 2023）在 24 个 ICS 控制二进制上实现 100\% 漏洞定位成功率和 $\sim$124--222 ms 补丁生成时间，运行时开销仅 0.05--17.4 $\mu$s，为工业控制系统提供零停机热补丁。在侧信道方面，CipherH~\cite{deng2023cipherh}（USENIX Security 2023）在 OpenSSL、MbedTLS、WolfSSL 中检测到 236 个密文侧信道漏洞点（精度 $\sim$96.2\%），发现 WolfSSL 中此前未知的漏洞并推动官方发布补丁。CacheQL~\cite{yuan2023cacheql}（USENIX Security 2023）使用 Shapley 值量化缓存侧信道泄露量（如 OpenSSL AES-128 泄露 128 比特），对 13 个常量时间函数正确报告零泄露，验证了精度。

在代码安全增强方面，DeGPT~\cite{degpt2024} 使用 LLM 优化反编译器输出，改善逆向工程中的代码可读性。SigmaDiff~\cite{sigmadiff2024} 通过语义感知深度图匹配实现伪代码差异分析。GNNic~\cite{gnnic2024} 利用抽象相似度查找"失散"的函数兄弟关系。

在 API 安全方面，Midas/ChatDetector~\cite{midas2025}（NDSS 2025）以 98.21\% 精度识别 165 个 RM-API 对，发现 115 个安全 bug，链式思维+交叉验证将精度从 77.97\% 提升至 98.21\%。API 参数安全规则生成~\cite{api_rules2025} 使用 LLM 生成 API 误用检测的安全规则。NodeMedic-FINE~\cite{song2025nodemedic} 对 Node.js 包进行细粒度漏洞检测和利用合成。YuraScanner~\cite{yurascanner2025}（NDSS 2025）利用 LLM 驱动任务导向的 Web 应用扫描，在 20 个 web 应用上发现比其他工具多 55.19\% 的表单和 35.16\% 的 URL，发现 12 个零日 XSS 漏洞。

\subsubsection{LLM 代码生成安全}
LLM 生成的代码本身可能引入安全缺陷，这使得审计不仅面向手写代码，也面向 AI 生成代码。Ullah 等~\cite{pearce2023copilot_correctness}（IEEE S\&P 2024）在 228 个代码场景、8 个 LLM 和 17 种提示技术上评估发现：改变函数名导致 PaLM2 在 26\% 情况下给出错误答案，添加库函数导致 GPT-4 在 17\% 情况下失败，且所有模型即使在温度=0.0 时仍存在非确定性输出。Niu 等~\cite{niu2023codexleaks}（USENIX Security 2023）发现约 8\% 的 prompt 可从 Codex 模型中泄露隐私信息（邮箱、姓名、SSN、API 密钥）。

INDICT~\cite{indict2024}（NeurIPS 2024）在代码生成中引入双重批评者内部对话机制，在 7B--70B 参数的 8 种语言上实现 +10\% 代码质量绝对提升，HarmBench 安全分类从 23.6\% 提升至 51.5\%。LeDex~\cite{ledex2024} 训练 LLM 更好地自调试和解释代码。WorldCoder~\cite{worldcoder2024} 通过构建世界模型和编写代码评估 LLM agent 能力。Verified Code Transpilation~\cite{verified_transpile2024} 使用 LLM 进行经过验证的代码转译。

REPOCOD~\cite{repocod2025}（ACL 2025）在 980 个仓库级 Python 任务上发现\textbf{没有 LLM 超过 30\% pass@1}（GPT-4o 最高仅 27.4\%），50.8\% 任务需要仓库级上下文。Can You Trust Code Copilot~\cite{trust_copilot2025}（ACL 2025）在 20 个 LLM 上发现：训练数据中 4.4\% 包含漏洞，与 LLM 生成代码的漏洞分布高度相关；Claude-3 漏洞修复率 66.25\%、GPT-4o 63.94\%，但大多数 LLM 检测 F1>80\% 却仍生成含漏洞代码。SecRepoBench~\cite{secrepobench2025}（LLM4Code 2026）在 318 个 C/C++ 仓库级补全任务上发现最强独立模型 GPT-5 仅 39.3\% secure-pass@1，代码 agent 显著提升正确性但安全性增益有限（仅 +4.7\%）。

这些工作共同说明：LLM 代码生成是审计的新攻击面，审计系统需要能够检测 AI 生成代码中的安全缺陷，且不能假设 AI 生成代码比手写代码更安全。

\subsection{Prompt 设计、安全感知微调与交互策略}
Prompt 与 fine-tuning 是审计性能的重要决定因素。同一模型在不同提示模板下输出差异显著，安全任务中少量措辞变化即可影响误报/漏报~\cite{icode_reviewer_2025,secure_reviewer_2025,eval_code_review_2025,liu_security_review_2024,pornprasit_2024_gpt35,zhang2024prompt}。Niu 等~\cite{niu2025largetomammoth} 在 NDSS 2025 的系统评估中明确发现：\textbf{prompt 设计比模型规模更重要}——7B 参数模型在精心 prompt 下可超越 405B 模型的零样本性能。

安全审计 prompt 不是让模型自由发挥，而是强约束式判定协议，应包含以下要素：
\begin{itemize}[leftmargin=*,noitemsep]
  \item \textbf{输出类别约束}：明确枚举 benign/malicious/uncertain，禁止模糊回答
  \item \textbf{证据优先级}：要求模型先列出证据再给结论，而非先结论后解释
  \item \textbf{Benign trap 豁免}：显式说明教学示例、测试用例、包装代码不应被标记为恶意
  \item \textbf{JSON 格式约束}：强制结构化输出，便于自动化解析和聚合
  \item \textbf{不确定性处理}：明确定义 UNCERTAIN 的触发条件（证据不足、多证据冲突）
\end{itemize}

在微调策略上，Mixture-of-prompts~\cite{icode_reviewer_2025}（华为，ASE 2025）将不同样本交给不同提示策略并结合静态分析路由，实现 63.98\% F1 和 84\% 生产接受率（C/C++）。Secure-aware fine-tuning~\cite{secure_reviewer_2025}（ICSE 2026）将安全判断偏好嵌入模型参数，结合 RAG 达到 F1=71.98\% 和 SecureBLEU=29.31。Chen 等~\cite{chen2026lmcode} 的大规模实证研究发现任务特定微调一致优于通用预训练模型。Pornprasit 和 Tantithamthavorn~\cite{pornprasit_2024_gpt35} 发现 few-shot learning 可将 GPT-3.5 性能提升至少 46\%，但 persona 提示反而降低性能。Effectiveness~\cite{effectiveness_vuln_llm_2025} 从实证评测角度说明模型即使能识别可疑模式，也未必能稳定完成漏洞成因推理——识别和推理是两个不同层次的能力。

Zhang 等~\cite{zhang2024prompt} 证明精心设计的 prompt 可显著改善 ChatGPT 零样本漏洞检测，但改善幅度高度依赖漏洞类型：CWE-79（XSS）等模式化漏洞改善明显，而 CWE-502（反序列化）等需要深层语义理解的漏洞改善有限。这进一步支持了本文的核心主张：prompt 优化是必要但不充分的条件，高可靠审计需要系统级架构而非单轮问答。

\subsection{Benchmark 设计与评估协议}
LLM 辅助审计的研究质量高度依赖 benchmark。2024--2025 年见证了 benchmark 设计的爆炸式增长。图~\ref{fig:benchmark_coverage} 给出主要 benchmark 对不同威胁维度的覆盖情况。

\begin{figure}[htbp]
\centering
\begin{tikzpicture}
\definecolor{filled}{RGB}{76,175,80}
\definecolor{empty}{RGB}{245,245,245}
\def\cw{1.8cm}
\def\rh{0.85cm}

% Column headers (rotated)
\foreach \i/\name [count=\ci from 0] in {
  0/SecLLMHolmes, 1/InjecAgent, 2/AgentDojo, 3/JailbreakBench,
  4/AgentSecBench, 5/WASP, 6/REPOCOD, 7/SecRepoBench} {
  \node[font=\scriptsize\bfseries,rotate=45,anchor=south west] at (\i*\cw+0.1*\cw, 6.2*\rh) {\name};
}

% Row headers
\foreach \j/\name in {0/漏洞检测, 1/间接提示注入, 2/越狱攻击, 3/Agent 行为, 4/模型后门, 5/代码生成安全} {
  \node[font=\small,anchor=east] at (-0.2, \j*\rh+0.5*\rh) {\name};
}

% Draw all cells first (empty)
\foreach \i in {0,...,7} {
  \foreach \j in {0,...,5} {
    \fill[empty,rounded corners=2pt] (\i*\cw, \j*\rh) rectangle ++(\cw, \rh);
    \draw[gray!20,rounded corners=2pt] (\i*\cw, \j*\rh) rectangle ++(\cw, \rh);
  }
}

% Fill covered cells
\def\fillcell#1#2{
  \fill[filled,rounded corners=2pt] (#1*\cw, #2*\rh) rectangle ++(\cw, \rh);
  \node[font=\large,text=white] at (#1*\cw+0.5*\cw, #2*\rh+0.5*\rh) {\cmark};
}

% Coverage data
\fillcell{0}{0}  % SecLLMHolmes: 漏洞检测
\fillcell{1}{1}  % InjecAgent: IPI
\fillcell{2}{1}  % AgentDojo: IPI
\fillcell{2}{3}  % AgentDojo: Agent行为
\fillcell{3}{2}  % JailbreakBench: 越狱
\fillcell{4}{1}  % AgentSecBench: IPI
\fillcell{4}{2}  % AgentSecBench: 越狱
\fillcell{4}{3}  % AgentSecBench: Agent行为
\fillcell{5}{3}  % WASP: Agent行为
\fillcell{6}{5}  % REPOCOD: 代码生成
\fillcell{7}{5}  % SecRepoBench: 代码生成

% Gap annotation
\draw[red!60!black,very thick,->,>=latex] (8.2*\cw, 4*\rh+0.5*\rh) -- (8.8*\cw, 4*\rh+0.5*\rh);
\node[font=\small\bfseries,text=red!60!black,anchor=west] at (8.9*\cw, 4*\rh+0.5*\rh) {空白};

% Summary row
\node[font=\scriptsize,text=gray!60,anchor=west] at (0, -0.5*\rh) {覆盖率：漏洞 1/8 · IPI 3/8 · 越狱 2/8 · Agent 3/8 · 后门 0/8 · 代码生成 2/8};

\end{tikzpicture}
\caption{主要 Benchmark 对威胁维度的覆盖矩阵}
\label{fig:benchmark_coverage}
\end{figure}

\subsubsection{漏洞检测 Benchmark}
LLM vs Static~\cite{llm_vs_static_2025} 将 LLM 与静态工具放在同一基准比较。Ullah 等~\cite{ullaha_vuln_reasoning_2024} 和 Paquette 等~\cite{paquette2024llmscannot} 分别从推理和识别两个角度建立了 LLM 漏洞检测评测框架。Zhang 等~\cite{zhang2024vulllmeval} 提出 VulnLLMEval 统一评估框架。Fair Security Eval~\cite{fair_security_eval_2025} 强调公平评估原则。Niu 等~\cite{niu2025largetomammoth} 首次对 7B--405B 规模模型进行系统对比。

\subsubsection{Agent 安全 Benchmark}
InjecAgent~\cite{zhan2024injecagent} 和 InjectBench~\cite{kong2024injectbench} 聚焦间接提示注入。Liu 等~\cite{liu2024promptinjbench} 提供统一形式化和综合 benchmark。AgentDojo~\cite{debenedetti2024agentdojo} 和 Agent Security Bench~\cite{zhang2024asb} 覆盖完整 agent 安全场景。WASP~\cite{evtimov2025wasp} 关注 Web agent 安全。Zhang 等~\cite{zhang2025tasklevel} 提供任务级越狱评估。

\subsubsection{越狱 Benchmark}
JailbreakBench~\cite{chao2024jailbreakbench}（NeurIPS 2024）建立标准化评测协议（100 个行为、4 种攻击基线、5 种防御基线），其 Prompt+RS 基线对 Vicuna/Llama-2/GPT-3.5/GPT-4 分别实现 89/90/93/78\% ASR，Erase-and-Check 是最强防御但 Prompt+RS 仍能达 8--25\%；GCG 需 25.6 万次查询和 1700 万 token，而 Prompt+RS 仅需 2--1K 次查询。Bag of Tricks~\cite{xu2024bagoftricks} 分解攻击组件，发现语义混淆和角色扮演主导字符级扰动。MasterKey~\cite{deng2024masterkey} 实现自动化越狱评估并发现越狱 prompt 的 GPT 专属性。PAPILLON~\cite{gong2025papillon} 和 TwinBreak~\cite{krauss2025twinbreak} 将攻击成功率推到 89--98\% 范围。LLM-Fuzzer~\cite{yu2024fuzzer} 将 fuzzing 方法论引入评测。

\subsubsection{代码审查 Benchmark}
CRUXEval~\cite{cruxeval_2024}（ICML 2024）在 800 个 Python 函数上评测代码执行推理能力：GPT-4 在输入预测上 pass@1=67.1\%、输出预测 63.4\%，链式思维提升至 75\%/81\%；但最佳开源模型仅 47.9\%/44\%，且链式思维对开源模型\textbf{无改善}——揭示闭源与开源模型在代码推理上的巨大鸿沟。Peng 等~\cite{peng_2025_icodereviewer} 报告 63.98\% F1 的安全审查基准。Cihan 等~\cite{cihan_2025_llm_review_eval} 在 492 个 AI 代码块上评测 LLM 审查能力。

\section{综合比较}

\begin{table}[htbp]
\caption{方法家族综合比较矩阵}
\label{tab:method_matrix}
\centering
\small
\begin{tabularx}{\textwidth}{@{} p{0.10\textwidth}X X X X X p{0.05\textwidth} @{}}
\rowcolor{tabhead}
\tabheader{家族} & \tabheader{核心目标} & \tabheader{证据输入} & \tabheader{主要优势} & \tabheader{主要短板} & \tabheader{本仓库样本} & \tabheader{可信度} \\
\rowcolor{tabalt} 直接代码审查 & 局部安全判断 & 函数、补丁、PR & 接入简单 & 上下文弱，18--68\% & B1-08 vs M1-01 & 高 \\
漏洞检测/修复 & 定位并修复 & 代码片段、上下文 & 任务边界清楚 & ML 学表面模式 & eval(label) vs eval() & 高 \\
\rowcolor{tabalt} 静态分析协同 & 召回+裁决 & 规则命中、调用链 & 可追溯、入 CI & SAST 76\% 不相关 & M3-02 vs B3-01 & 高 \\
结构化/RAG & 高信噪比证据 & CPG/CFG、切片 & 降噪、补上下文 & 前处理和检索污染 & M5-01 & 中--高 \\
\rowcolor{tabalt} Agent 审计 & 运行时行为 & prompt、工具、记忆 & 覆盖真实攻击面 & 能力越强越脆弱, 防御降 15--20\% & 行为链类比 & 中 \\
越狱攻防 & 鲁棒性评估 & 对抗 prompt & 系统化评测 & 89--98\% ASR, 理论不可避免 & N/A & 中 \\
\rowcolor{tabalt} 模型完整性 & 信任链验证 & 权重、激活 & 确保审计器可信 & 后门化→全失效 & N/A & 中 \\
LLM 安全测试 & 主动漏洞发现 & LLM fuzzing & 扩大覆盖面 & 依赖 LLM 能力 & N/A & 中 \\
\rowcolor{tabalt} 代码生成安全 & AI 代码审计 & 生成代码 & 覆盖新攻击面 & AI 代码引入缺陷 & N/A & 中 \\
Prompt/微调 & 稳定判定 & 约束 prompt & 改善一致性 & 版本标签敏感 & B3-01/B3-04 & 中 \\
\rowcolor{tabalt} Benchmark & 评估有效性 & manifest、样本 & 可复跑比较 & 来源泄漏风险 & audit-100 & 高 \\
\end{tabularx}
\end{table}

\subsection{审计器模型与使用方式}
本文严格区分被审计代码包与审计器模型。\path{LogisticRegression}、\path{TEE-test}、\path{Retina-DKD} 是 subject-side 样本来源；Claude、GPT、Gemini、Qwen、DeepSeek、StarCoder 等才是 auditor-side 模型候选。表~\ref{tab:model_matrix} 给出审计器类型比较。

\begin{table}[htbp]
\caption{审计器模型与使用方式对比}
\label{tab:model_matrix}
\centering
\small
\begin{tabularx}{\textwidth}{@{} p{0.10\textwidth}X X X X p{0.08\textwidth} @{}}
\rowcolor{tabhead}
\tabheader{类型} & \tabheader{代表模型} & \tabheader{推荐角色} & \tabheader{优势} & \tabheader{风险} & \tabheader{状态} \\
\rowcolor{tabalt} 前沿闭源 & Claude、GPT-4.1、Gemini & 语义裁决、边界仲裁 & 推理强，跨文件 & 成本、版本漂移 & 候选 \\
本地代码 & Qwen2.5-Coder、DeepSeek & 可复现批量审计 & 成本低，可离线 & 格式不稳 & 已接入 \\
\rowcolor{tabalt} 静态规则 & Python/Go 规则 & 高召回、fail-closed & 可解释、可版本化 & 易被变形绕过 & 满分 \\
混合系统 & 静态+结构化+LLM+策略 & 生产级审计闭环 & 兼顾召回与追溯 & 工程最复杂 & 目标 \\
\end{tabularx}
\end{table}

\subsection{当前评测结果与含义}
表~\ref{tab:result_matrix} 汇总本仓库关键评测阶段。

\begin{table}[htbp]
\caption{当前仓库评测阶段对比}
\label{tab:result_matrix}
\centering
\small
\begin{tabularx}{\textwidth}{@{} p{0.10\textwidth}X X X X @{}}
\rowcolor{tabhead}
\tabheader{阶段} & \tabheader{样本语义} & \tabheader{代表结果} & \tabheader{指标含义} & \tabheader{后续动作} \\
\rowcolor{tabalt} 初始静态 & 来源未隔离 & TP=43, FP=22, TN=28, FN=7 & 有召回但误报高 & 修正来源隔离 \\
Ollama 后 & LLM 不稳定 & TP=41, FP=22, TN=28, FN=9, 率=0.57 & 可运行≠可靠 & 加强 JSON 约束 \\
\rowcolor{tabalt} 来源隔离 & b/m 分离 & TP=50, FP=0, TN=50, FN=0 & 可解释但偏易 & 增加跨文件化 \\
跨文件化 & 文件+wrapper & 仍全分离 & 规则可全分离 & 提升隐蔽度 \\
\rowcolor{tabalt} difficulty6 & 100样本静态分离 & TP=50, FP=0, TN=50, FN=0 & 不能证明LLM增益 & 继续提高难度 \\
\end{tabularx}
\end{table}

\subsection{代表性 Benchmark 对比}
表~\ref{tab:bench_compare} 对比近年主要 benchmark，揭示领域评测覆盖面的结构性特征。

\begin{table}[t]
\caption{代表性 Benchmark 对比}
\label{tab:bench_compare}
\centering
\begin{tabularx}{\textwidth}{@{} p{0.12\textwidth}p{0.08\textwidth}X X X X @{}}
\rowcolor{tabhead}
\tabheader{Benchmark} & \tabheader{Venue} & \tabheader{任务} & \tabheader{规模} & \tabheader{核心发现} & \tabheader{审计启示} \\
\rowcolor{tabalt} SecLLMHolmes~\cite{paquette2024llmscannot} & S\&P'24 & 漏洞识别+推理 & 多 CWE 多模型 & LLM 无法可靠区分漏洞 & 不应单靠 LLM \\
InjecAgent~\cite{zhan2024injecagent} & ACL'24 & 间接提示注入 & 1,054 用例, 30 agents & GPT-4 ASR 24--47\%; 微调降至 3.8\% & 微调优于 prompt \\
\rowcolor{tabalt} AgentDojo~\cite{debenedetti2024agentdojo} & NeurIPS'24 & Agent 安全 & 97 任务, 629 安全测试 & 工具过滤$\to$7.5\% ASR; 效用损失 15--20\% & 安全-效用权衡 \\
JailbreakBench~\cite{chao2024jailbreakbench} & NeurIPS'24 & 越狱评测 & 100 行为, 4 攻击 5 防御 & Prompt+RS ASR 78--93\% & 防御效果微弱 \\
\rowcolor{tabalt} Agent Sec Bench~\cite{zhang2024asb} & ICLR'25 & Agent 综合 & 10 场景, 13 LLM, 400+ 工具 & 平均 ASR=84.3\%; 防御有效性 46.9\% & 能力越强越脆弱 \\
WASP~\cite{evtimov2025wasp} & NeurIPS'25 & Web Agent & 84 对抗任务 & 劫持 85.7\% 但目标完成 $\leq$17\% & "安全源于无能" \\
\rowcolor{tabalt} REPOCOD~\cite{repocod2025} & ACL'25 & 仓库级生成 & 仓库级任务 & LLM 尚不能替代程序员 & 覆盖 AI 生成代码 \\
SecRepoBench~\cite{secrepobench2025} & arXiv'25 & 安全补全 & 真实仓库 & 安全补全能力有限 & 系统需独立验证 \\
\end{tabularx}
\end{table}

\subsection{相关综述对比}
表~\ref{tab:survey_compare} 将本文与近期同类综述进行对比。

\begin{table}[t]
\caption{相关综述对比}
\label{tab:survey_compare}
\centering
\begin{tabularx}{\textwidth}{@{} p{0.12\textwidth}p{0.08\textwidth}X X X @{}}
\rowcolor{tabhead}
\tabheader{综述} & \tabheader{Venue} & \tabheader{覆盖范围} & \tabheader{本文差异} & \tabheader{互补价值} \\
\rowcolor{tabalt} Sheng~\cite{sheng2026llmsecurity} & CSUR'26 & 80+ 篇, 33 LLM & 覆盖完整审计链 & GPT-4 最常用（29 次） \\
Chen~\cite{chen2026lmcode} & TOSEM'26 & 68 篇, 攻击/防御 & 增加 IPI、越狱、后门 & 后门+对抗分类学 \\
\rowcolor{tabalt} Kaniewski~\cite{kaniewski2026detecting} & TOSEM'26 & 263 篇（最大 SLR） & 覆盖 agent 和审查 & 碎片化, 缺标准化 \\
Zhou~\cite{zhou2025vulndetrepair} & TOSEM'25 & 58 篇, 37 LLM & 增加安全测试和治理 & 最优检测 67.6\%, 修复 $\sim$20\% \\
\rowcolor{tabalt} Xu~\cite{xu2026duality} & AI Rev'26 & Agent 双重性 & 聚焦审计应用 & 引用双重性框架 \\
SoK AVR~\cite{sok2025avr} & SoK'25 & 自动修复 & 审计视角审视可信性 & 引用修复分类学 \\
\rowcolor{tabalt} Yang~\cite{yang_2026_roadmap} & arXiv'26 & 审查路线图 & 聚焦安全审计 & 引用范式转变 \\
\end{tabularx}
\end{table}

\subsection{跨方法交叉分析}
综合 158 条引用，本文识别出三个贯穿所有方法家族的交叉主题：

\textbf{主题一：语义能力 $\neq$ 安全可靠性}。
代码审查（§III-A）中活跃用户研究接受率仅 8.1\%~\cite{olewicki_2024_llm_review_impact}，训练数据约 25\% 为噪声~\cite{tufano_2024_strengths_weaknesses}；漏洞检测（§III-B）中 ML4VD 模型准确率低于随机猜测~\cite{risse2024limits}，更大模型反而不如小模型（LLaMA-3 70B F1=5.10\% vs Gemma 7B F1=57.98\%）~\cite{niu2025largetomammoth}；agent 安全（§III-E）中平均 ASR 达 84.30\% 且能力越强越脆弱~\cite{zhang2024asb}；越狱（§III-F）中 Prompt+RS 对 4 个 LLM 实现 78--93\% ASR~\cite{chao2024jailbreakbench}。五个独立领域的实证数据共同指向同一结论：LLM 的语义理解能力并不自动转化为安全判断的可靠性。这意味着审计系统必须将 LLM 作为组件之一，而非唯一裁决器。

\textbf{主题二：评估基础设施决定研究可信度}。
Kaniewski 等~\cite{kaniewski2026detecting} 在 263 篇研究中发现领域高度碎片化——各研究使用自定义数据集、指标和数据划分，难以跨研究比较。从代码审查数据约 25\% 噪声~\cite{tufano_2024_strengths_weaknesses}到越狱防御在标准化测试下效果微弱~\cite{chao2024jailbreakbench}，从 ML4VD 数据增强仅恢复 2.5--5.9\% 准确率~\cite{risse2024limits}到所有 agent 防御均导致 15--20\% 效用损失~\cite{debenedetti2024agentdojo}，评估方法论本身成为制约研究进展的瓶颈。图~\ref{fig:benchmark_coverage} 的覆盖热力图进一步揭示：模型后门维度至今无专门 benchmark，这是一个亟待填补的空白。

\textbf{主题三：攻击面持续扩大}。
从静态代码漏洞（§III-B）到 agent 执行链（§III-E），从 prompt injection（InjecAgent 显示 GPT-4 ASR 24--47\%~\cite{zhan2024injecagent}）到模型后门（InstrBackdoor 在所有 LLM 上 ASR=1.0~\cite{zhang2024instrbackdoor}），从手写代码到 AI 生成代码~\cite{pearce2023copilot_correctness}，审计攻击面在 2013--2026 年间持续扩大（图~\ref{fig:timeline}）。AgentFuzz 在 20 个 agent 应用中发现 34 个 0-day~\cite{liu2025taintstyle}，98,380 个社区技能中存在大量恶意模式~\cite{liu2026malskills}，KV-cache 共享导致跨租户泄露~\cite{wu2025promptleakage}。每一次技术演进都引入新的攻击向量：RAG 引入检索投毒（PRA-RAG 将 ASR 从 46\% 降至 1\%~\cite{tan2026prarag}），agent 引入工具滥用，微调引入后门植入。审计系统必须设计为可演进架构，能够适应不断变化的威胁格局。

\section{威胁有效性与局限}

\subsection{文献选择偏差}
本文采用结构化叙述方法，存在英文文献偏置、arXiv 预印本偏置、新近研究曝光度偏置，以及关键词筛选偏差。本文不声称穷尽所有 LLM 安全审计论文。

\subsection{概念不一致与基准异构}
不同论文对漏洞、安全缺陷、风险、恶意、可疑和高危的定义并不一致。模型、数据集、上下文粒度、prompt、温度、后处理与标注标准也可能不同。Risse 和 Böhme~\cite{risse2024limits} 明确展示了这种异构性：顶级 ML4VD 模型在语义等价代码上表现不一致，说明基准本身的定义就可能导致误导性结论。

\subsection{时间漂移与版本漂移}
LLM 研究迭代极快。本文曾出现 \risk{llm\_available\_rate}=0 的情况。Kaniewski 等~\cite{kaniewski2026detecting} 在工业评估中发现模型版本和 prompt 变化显著影响结果。因此，研究必须记录 \risk{rule\_set\_version}、\risk{prompt\_version}、\risk{auditor\_version}、运行时间和失败状态。

\subsection{攻防持续博弈}
越狱攻防呈现持续升级态势。Su 等~\cite{su2024missionimpossible} 从理论上证明越狱在统计意义上不可避免。攻击端从 MasterKey（21.58\% 成功率）~\cite{deng2024masterkey} 演进到 PAPILLON（>90\% ASR，困惑度 34.61 绕过检测）~\cite{gong2025papillon} 和 TwinBreak（89--98\% ASR，<5 分钟）~\cite{krauss2025twinbreak}。防御端从 Spotlighting（ASR→0\%）~\cite{hines2024spotlighting} 到 StruQ（TAP 97\%→9\%）~\cite{chen2025struq} 再到 InstructDetector（仅需 200 样本达 99.6\% 检测率）~\cite{wen2025instrdetect}。Gong 等~\cite{gong2025safetyMisalign} 揭示对齐技术产生"对齐税"（HarmBench ASR 42\%→36.95\%）。这意味着安全审计系统不能仅依赖模型自身的安全对齐，而必须在系统层面实施防御纵深。

\subsection{综述自身边界}
本文没有采用完整 PRISMA 流程，也未对所有候选论文做量化元分析。因此本文定位是结构化、证据驱动的学术综述。若后续升级为正式系统综述，需要补充检索式全文、数据库命中数量、去重规则、筛选流程图、双人标注一致性、质量评估量表、排除论文清单和完整引用。

表~\ref{tab:threats_summary} 汇总上述威胁及其缓解策略。

\begin{table}[htbp]
\caption{威胁有效性总结}
\label{tab:threats_summary}
\centering
\small
\begin{tabularx}{\linewidth}{@{} p{0.20\linewidth}X p{0.20\linewidth} @{}}
\rowcolor{tabhead}
\tabheader{威胁类型} & \tabheader{具体表现} & \tabheader{缓解策略} \\
\rowcolor{tabalt} 文献选择偏差 & 英文/arXiv/新近偏置 & 多源交叉验证，标注证据等级 \\
概念异构 & 漏洞/风险/可疑定义不一 & 统一术语，引用形式化定义 \\
\rowcolor{tabalt} 版本漂移 & 模型/prompt 变化影响结果 & 固定版本号，记录失败状态 \\
攻防博弈 & 越狱不可避免，对齐税 & 系统级纵深防御，非单点对齐 \\
\rowcolor{tabalt} 综述边界 & 非 PRISMA，无元分析 & 定位为结构化综述，标注局限 \\
\end{tabularx}
\end{table}

\section{研究议程}

基于本文 158 条引用的分析和 100 样本 benchmark 的迭代经验，提出以下六个研究方向。每个方向给出具体问题、量化目标和验证方法。

\subsection{从单点分类走向证据链审计}
\textbf{问题}：当前 LLM 审计多为单轮问答，缺乏证据链追溯能力。

\textbf{方向}：未来的 LLM 审计应从"这段代码是否危险"转向"哪些证据能支撑危险判断"。代码、规则、调用链、检索记录和输出解释应共享统一表示，并能回溯到具体文件、行号、证据来源与版本。Mirsky 等~\cite{mirsky2023vulchecker} 的行级定位和 Kim 等~\cite{kim2025san2patch} 的 sanitizer 日志驱动修复代表了这一方向。

\textbf{量化目标}：每个审计结论应关联至少 3 类独立证据（规则命中、调用链、历史案例）；证据链缺失时应自动升级为 UNCERTAIN。

\textbf{验证方法}：构建包含证据链标注的 benchmark，测量单证据 vs 多证据判断的 F1 差异。

\subsection{从单模型走向多阶段系统}
\textbf{问题}：单模型审计无法兼顾召回、精度和可解释性。

\textbf{方向}：高可靠系统应采用多阶段架构。Wu 等~\cite{wu2025isolateGpt} 的执行隔离和 Chen 等~\cite{chen2025struq} 的结构化查询提供了系统级防御的参考。不确定时应 fail-closed。

\textbf{量化目标}：第一阶段静态分析召回率 $\geq 95\%$，第二阶段 LLM 裁决将误报降低 $\geq 60\%$，第三阶段策略层确保 \risk{llm\_unavailable} 时 100\% fail-closed。

\textbf{验证方法}：在 100 样本 benchmark 上测量各阶段独立指标和端到端指标，对比单阶段 vs 多阶段的 precision-recall 曲线。

图~\ref{fig:pipeline} 展示了推荐的多阶段审计流水线架构。

\begin{figure}[htbp]
\centering
\begin{tikzpicture}[
  stage/.style={draw=#1!60!black,fill=#1!12,rectangle,rounded corners=5pt,minimum width=2.8cm,minimum height=1.4cm,font=\small\bfseries,text=#1!60!black,align=center,drop shadow},
  metric/.style={font=\scriptsize\itshape,text=#1!60!black,align=center},
  data/.style={draw=gray!40,fill=gray!5,ellipse,minimum width=2cm,minimum height=0.9cm,font=\footnotesize,align=center,text=black!70},
  decision/.style={draw=orange!60!black,fill=orange!10,diamond,aspect=2.2,minimum width=1.8cm,minimum height=1.1cm,font=\footnotesize\bfseries,text=orange!60!black,align=center},
  output/.style={draw=#1!50!black,fill=#1!10,rectangle,rounded corners=3pt,minimum width=1.6cm,minimum height=0.6cm,font=\footnotesize,text=#1!60!black,align=center},
  arrow/.style={->,>=latex,line width=1.2pt,color=black!50}
]

% Input
\node[data] (input) at (-2.5,0) {代码/配置\\依赖/PR};

% Four stages
\node[stage=blue] (s1) at (1,0) {阶段 1\\静态分析};
\node[stage=orange] (s2) at (5,0) {阶段 2\\结构化增强};
\node[stage=purple] (s3) at (9,0) {阶段 3\\LLM 语义裁决};
\node[stage=teal] (s4) at (12.5,0) {阶段 4\\策略层};

% Metrics below stages
\node[metric=blue] at (1,-1.2) {召回 $\geq 95\%$};
\node[metric=orange] at (5,-1.2) {信噪比提升};
\node[metric=purple] at (9,-1.2) {误报降 $\geq 60\%$};
\node[metric=teal] at (12.5,-1.2) {100\% fail-closed};

% Decision diamond
\node[decision] (dec) at (12.5,-2.8) {阻断?};

% Outputs
\node[output=red] (block) at (15,-2.3) {阻断};
\node[output=green!70!black] (pass) at (12.5,-4.2) {放行};
\node[output=yellow!70!black] (review) at (10,-2.3) {人工复核};

% Main flow arrows
\draw[arrow] (input) -- (s1);
\draw[arrow] (s1) -- (s2);
\draw[arrow] (s2) -- (s3);
\draw[arrow] (s3) -- (s4);
\draw[arrow] (s4) -- (dec);

% Decision outputs
\draw[arrow] (dec) -- node[above,font=\scriptsize] {高危} (block);
\draw[arrow] (dec) -- node[right,font=\scriptsize] {安全} (pass);
\draw[arrow] (dec) -- node[above,font=\scriptsize] {不确定} (review);

% Feedback loops (dashed)
\draw[arrow,dashed,color=purple!50] (s3.north) -- ++(0,0.7) -| node[pos=0.25,above,font=\scriptsize,text=purple!60!black] {证据不足} (s2.north);
\draw[arrow,dashed,color=teal!50] (s4.north) -- ++(0,1.1) -| node[pos=0.25,above,font=\scriptsize,text=teal!60!black] {LLM 不可用 $\to$ 保守} (s1.north);

% Data annotations below stages
\node[font=\tiny,text=gray!60,below=1.8cm of s1] {规则命中、污点路径};
\node[font=\tiny,text=gray!60,below=1.8cm of s2] {CPG/CFG、切片};
\node[font=\tiny,text=gray!60,below=1.8cm of s3] {语义解释、评分};
\node[font=\tiny,text=gray!60,below=1.8cm of s4] {策略判定、日志};

\end{tikzpicture}
\caption{推荐的多阶段审计流水线架构}
\label{fig:pipeline}
\end{figure}

\subsection{从漏洞检测走向应用行为审计}
\textbf{问题}：静态代码审计无法覆盖运行时威胁（工具调用、提示注入、记忆污染）。

\textbf{方向}：未来真正危险的部分常常不在静态代码，而在运行时行为：工具调用、提示注入、记忆污染、检索污染和多轮规划。Liu 等~\cite{liu2025taintstyle} 的 agent 污点分析和 Liu 等~\cite{liu2026malskills} 的恶意技能检测代表了行为链审计的前沿。

\textbf{量化目标}：行为审计应覆盖至少 5 类运行时威胁（工具滥用、IPI、记忆篡改、检索投毒、多轮偏离），每类威胁检测率 $\geq 80\%$。

\textbf{验证方法}：构建包含行为链标注的 agent benchmark，测量静态审计 vs 行为审计的威胁覆盖率差异。

\subsection{从准确率导向走向可治理性导向}
\textbf{问题}：高准确率不等于高可信度。LLM 版本漂移、prompt 变化和数据污染使指标不可复现。

\textbf{方向}：高风险审计更需要可复跑、可解释、可追溯、可版本化、失败显式化和不确定性可管理。JailbreakBench~\cite{chao2024jailbreakbench} 的标准化评测协议和 SoK~\cite{sok2025avr} 的评估分类学提供了 benchmark 治理的参考。

\textbf{量化目标}：每次评测必须记录 \risk{rule\_set\_version}、\risk{prompt\_version}、\risk{auditor\_version}、\risk{llm\_available\_rate}、\risk{fail\_closed\_count}；结果包应包含 confusion matrix、sample-level 结果和失败样本分析。

\textbf{验证方法}：同一 benchmark 在不同版本、不同时间、不同模型上复跑，测量指标漂移 $\leq 5\%$。

\subsection{从一般微调走向安全感知微调}
\textbf{问题}：通用预训练模型在安全任务上表现不稳定。

\textbf{方向}：Chen 等~\cite{chen2026lmcode} 证明任务特定微调一致优于通用预训练。Peng 等~\cite{peng_2025_icodereviewer} 的混合提示和 SecureReviewer~\cite{secure_reviewer_2025} 的安全感知微调代表了两个互补方向。

\textbf{量化目标}：安全感知微调应在 benign trap 豁免上达到 100\% 准确率，在恶意样本上达到 $\geq 90\%$ F1，在 UNCERTAIN 判断上达到 $\geq 85\%$ 校准精度。

\textbf{验证方法}：在 100 样本 benchmark 上对比通用模型 vs 微调模型的 benign/malicious/uncertain 三分类指标。

\subsection{从单一 benchmark 走向 benchmark governance}
\textbf{问题}：当前 benchmark 易受来源泄漏、标签噪声和评估偏差影响。

\textbf{方向}：Benchmark 设计必须治理化：来源隔离、版本固定、标签说明、错误分解和结果归档都应作为研究的一部分。Agent Security Bench~\cite{zhang2024asb} 和 JailbreakBench~\cite{chao2024jailbreakbench} 的标准化方法论值得借鉴。

\textbf{量化目标}：每个 benchmark 应包含 manifest（样本元数据）、sample-results（逐样本结果）、confusion-matrix（分类统计）和 iteration-log（优化历史）；来源隔离应通过 hash 校验确保。

\textbf{验证方法}：benchmark 在开源后 6 个月内被至少 3 个独立团队复现，指标偏差 $\leq 10\%$。

\subsection{实践建议汇总}
表~\ref{tab:practice_recommendations} 汇总面向工程实践者的具体建议。

\begin{table}[htbp]
\caption{工程实践建议}
\label{tab:practice_recommendations}
\centering
\small
\begin{tabularx}{\linewidth}{@{} p{0.20\linewidth}X @{}}
\rowcolor{tabhead}
\tabheader{场景} & \tabheader{建议} \\
\rowcolor{tabalt} 代码审查自动化 & 接受率 $<50\%$ 时不应全自动部署；需人工复核边界样本 \\
漏洞检测 & 不依赖单一 LLM；静态+LLM 混合，SAST 负责召回，LLM 过滤误报 \\
\rowcolor{tabalt} Agent 安全审计 & 覆盖工具调用、外部输入和记忆；实施执行隔离~\cite{wu2025isolateGpt} \\
模型选型 & 本地模型做可复现基线，闭源模型做边界仲裁；prompt $>$ 规模~\cite{niu2025largetomammoth} \\
\rowcolor{tabalt} 不确定性处理 & UNCERTAIN 显式记录并升级；LLM 不可用时必须 fail-closed \\
Benchmark 使用 & 检查来源隔离、版本固定和标签说明；不信任未公开 manifest \\
\end{tabularx}
\end{table}

\section{结论}
LLM 已成为应用程序审计的重要组成部分，但最成熟的范式不是纯 LLM 替代传统审计，而是 LLM 与静态分析、结构化证据、检索机制、prompt 约束和策略控制协同工作。本文基于 200 余篇文献、158 条引用的综合分析显示：

\textbf{代码审查方面}，从 2013 年 Bacchelli 的基础研究到 2026 年的 AI 辅助路线图，MCR 经历了从人工到 LLM 辅助的完整演进。LLM 审查接受率约 18--84\%~\cite{olewicki_2024_llm_review_impact,peng_2025_icodereviewer}，数据质量和噪声是关键瓶颈~\cite{liu_2025_noisy_data,tufano_2024_strengths_weaknesses}，SAST 工具虽能覆盖 52\% 的漏洞变更但 76\% 警告不相关~\cite{charoenwet_2024_sast_secure_review}。

\textbf{漏洞检测方面}，从 VUDDY（2017）到 LLMxCPG（2025）的技术谱系展示了从规则到深度学习的演进。但 IEEE S\&P 和 USENIX Security 多项研究明确证明 LLM 无法可靠识别安全漏洞~\cite{paquette2024llmscannot,ullaha_vuln_reasoning_2024}，ML 模型学习表面模式~\cite{risse2024limits}，更大模型不优于小模型~\cite{niu2025largetomammoth}。自动修复取得显著进展——PatchAgent 达 92.13\% 修复成功率~\cite{yu2025patchagent}，SAN2PATCH 达 79.5\%~\cite{kim2025san2patch}，AppAtch 零日 F1 达 36.46\%~\cite{wen2025appatch}——但检测质量约束修复质量~\cite{sok2025avr}。

\textbf{Agent 安全方面}，间接提示注入~\cite{zhan2024injecagent,liu2024promptinjbench}和越狱攻击~\cite{deng2024masterkey,gong2025papillon}构成系统性威胁：AgentFuzz 在 20 个 agent 应用中发现 34 个 0-day（精度 100\%）~\cite{liu2025taintstyle}，WASP 显示 web agent 劫持率高达 85.7\%~\cite{evtimov2025wasp}，而防御降低效用的同时最多降低 30\% 攻击成功率~\cite{zhang2024asb}，越狱在统计意义上不可避免~\cite{su2024missionimpossible}。主动防御方面，IsolateGPT 的 hub-and-spoke 架构保持 100\% 功能完整性~\cite{wu2025isolateGpt}，CHeaT 框架在 CTF 上实现 100\% 防御成功率~\cite{ayzenshteyn2025cloak}。模型后门~\cite{wei2024lmsanitator,zhang2024instrbackdoor}和 AI 生成代码安全~\cite{pearce2023copilot_correctness,trust_copilot2025}进一步扩展了审计攻击面。

\textbf{未来方向}是构建证据链审计系统：静态分析产生候选（LLM F1=0.75--0.80 但定位有偏差~\cite{gnnic2024}），结构化增强提取证据（CPG/切片降噪），LLM 进行语义判定（需 fail-closed 处理 UNCERTAIN），RAG 补充 API 语义（需 PRA-RAG 级鲁棒聚合~\cite{tan2026prarag}），策略层决定阻断或放行，模型完整性验证纳入信任链（Beatrix F1=91.1\%~\cite{li2023beatrix}，LMSanitator 92.8\% 检测率~\cite{wei2024lmsanitator}），人工复核处理边界样本。Agent 审计需覆盖工具调用、IPI、记忆污染和检索投毒的完整行为链，且微调模型比 prompt 模型鲁棒 6$\times$ 的事实~\cite{zhan2024injecagent}提示审计器本身应优先使用微调而非零样本。

\begin{center}
\fbox{\parbox{0.9\linewidth}{\small
\textbf{核心主张（Take-home Message）}：
\begin{enumerate}[leftmargin=*,noitemsep]
  \item \textbf{证据链 $\succ$ 印象判断}：每个审计结论必须可追溯到具体文件、行号、规则和版本
  \item \textbf{多阶段系统 $\succ$ 单轮问答}：静态召回 $\to$ 结构化增强 $\to$ LLM 语义裁决 $\to$ 策略层阻断
  \item \textbf{行为链审计 $\succ$ 仅代码审查}：覆盖工具调用、IPI、记忆污染、检索投毒和多轮规划
  \item \textbf{模型完整性 $\succ$ 模型能力假设}：后门化模型使所有审计结论失效
  \item \textbf{失败显式化 $\succ$ 默认放行}：LLM 不可用时必须 fail-closed，不确定时升级为 UNCERTAIN
\end{enumerate}
}}
\end{center}

\appendix
\section{结论--证据追踪矩阵}
\begin{table}[t]
\caption{关键结论的证据追踪}
\label{tab:evidence_trace}
\centering
\begin{tabularx}{\textwidth}{@{} X X X p{0.06\textwidth} p{0.05\textwidth} @{}}
\rowcolor{tabhead}
\tabheader{关键结论} & \tabheader{文献证据} & \tabheader{样本/工件证据} & \tabheader{位置} & \tabheader{可信度} \\
\rowcolor{tabalt} LLM 无法可靠识别安全漏洞 & \cite{paquette2024llmscannot,ullaha_vuln_reasoning_2024,risse2024limits,niu2025largetomammoth} & ML4VD <50\% 准确率; LLaMA-3 F1=5.1\% & III-B & \textbf{高} \\
单轮审查受上下文影响 & \cite{liu_security_review_2024,cihan_2025_llm_review_eval,tufano_2024_strengths_weaknesses} & B1-08 vs M1-01 & III-A & \textbf{高} \\
\rowcolor{tabalt} 静态+LLM 适合工程审计 & \cite{program_analysis_survey_2025,llmxcpg_2025,charoenwet_2024_sast_secure_review} & M3-02 vs B3-01 & III-C & \textbf{高} \\
RAG 价值取决于证据质量 & \cite{tan2026prarag,hong_2025_rag_review,llmxcpg_2025} & M5-01 & III-D & 中--高 \\
\rowcolor{tabalt} Agent 审计须覆盖执行链 & \cite{debenedetti2024agentdojo,zhang2024asb,liu2025taintstyle,liu2026malskills} & ASR=84.3\%; 34 个 0-day & III-E & 中 \\
越狱理论上不可避免 & \cite{su2024missionimpossible,gong2025safetyMisalign,chao2024jailbreakbench,gong2025papillon} & PAPILLON >90\% ASR & III-F & 中 \\
\rowcolor{tabalt} 模型完整性是审计前提 & \cite{wei2024lmsanitator,zhang2024instrbackdoor,philosophers2025,yan2024llmbackdoor} & N/A & III-E.4 & 中 \\
AI 代码是新攻击面 & \cite{pearce2023copilot_correctness,indict2024,trust_copilot2025} & N/A & III-E.6 & 中 \\
\rowcolor{tabalt} Benchmark 治理决定可信性 & \cite{chao2024jailbreakbench,zhang2024asb,sok2025avr} & audit-100 满分 & III-G & \textbf{高} \\
\end{tabularx}
\end{table}

\begin{thebibliography}{99}

% === 原有 23 条 ===

\bibitem{liu_security_review_2024}
Y. Liu et al., ``An Insight into Security Code Review with LLMs: Capabilities, Obstacles and Influential Factors,'' arXiv:2401.16310, 2024.

\bibitem{zhou_vuln_repair_2024}
X. Zhou, S. Cao, X. Sun, and D. Lo, ``Large Language Model for Vulnerability Detection and Repair: Literature Review and the Road Ahead,'' ACM TOSEM, 2025.

\bibitem{llm_security_survey_2025}
``LLMs in Software Security: A Survey of Vulnerability Detection Techniques and Insights,'' arXiv:2502.07049, 2025.

\bibitem{slr_vuln_llm_2025}
``A Systematic Literature Review on Detecting Software Vulnerabilities with Large Language Models,'' arXiv:2507.22659, 2025.

\bibitem{program_analysis_survey_2025}
``A Contemporary Survey of Large Language Model Assisted Program Analysis,'' arXiv:2502.18474, 2025.

\bibitem{agent_audit_2026}
``Agent Audit: A Security Analysis System for LLM Agent Applications,'' arXiv:2603.22853, 2026.

\bibitem{rag_agent_security_2026}
``End-to-end Security Threats and Defenses in Retrieval-Augmented LLM Agents,'' Springer, 2026.

\bibitem{agent_security_duality_2026}
``LLM Agents Security Duality: A Comprehensive Survey,'' Artificial Intelligence Review, 2026.

\bibitem{malicious_skills_2026}
``Do Not Mention This to the User: Detecting and Understanding Malicious Agent Skills in the Wild,'' USENIX Security, 2026.

\bibitem{eval_code_review_2025}
``Evaluating Large Language Models for Code Review,'' arXiv:2505.20206, 2025.

\bibitem{auto_code_review_2024}
``Automated Code Review in Practice,'' arXiv:2412.18531, 2024.

\bibitem{overcorrection_2026}
``Are LLMs Reliable Code Reviewers? Systematic Overcorrection in Requirement Conformance Judgement,'' arXiv:2603.00539, 2026.

\bibitem{effectiveness_vuln_llm_2025}
``Understanding the Effectiveness of Large Language Models in Detecting Security Vulnerabilities,'' IEEE Xplore, 2025.

\bibitem{ullaha_vuln_reasoning_2024}
M. Ullah et al., ``LLMs Cannot Reliably Identify and Reason About Security Vulnerabilities,'' IEEE S\&P, 2024.

\bibitem{synergizing_static_llm_2025}
``Synergizing Static Analysis with Large Language Models for Vulnerability Discovery,'' arXiv:2509.15433, 2025.

\bibitem{llmxcpg_2025}
``LLMxCPG: Context-Aware Vulnerability Detection Through Code Property Graph-Guided LLMs,'' USENIX Security, 2025.

\bibitem{prompt_injection_agents_2025}
``Securing AI Agents Against Prompt Injection Attacks,'' arXiv:2511.15759, 2025.

\bibitem{llm_vs_static_2025}
``Large Language Models Versus Static Code Analysis Tools: A Systematic Benchmark,'' arXiv:2508.04448, 2025.

\bibitem{code_security_slr_2024}
``Large Language Models and Code Security: A Systematic Literature Review,'' arXiv:2412.15004, 2024.

\bibitem{lm_code_security_slr_2024}
``Security of Language Models for Code: A Systematic Literature Review,'' arXiv:2410.15631, 2024.

\bibitem{secure_reviewer_2025}
``SecureReviewer: Enhancing LLMs for Secure Code Review through Secure-aware Fine-tuning,'' arXiv:2510.26457, 2025.

\bibitem{icode_reviewer_2025}
``iCodeReviewer: Improving Secure Code Review with Mixture of Prompts,'' ASE, 2025.

\bibitem{fair_security_eval_2025}
``Automatic Security-Flaw Detection -- Towards a Fair Evaluation and Comparison,'' SoSyM, 2025.

% === 代码审查新增 ===

\bibitem{bacchelli_bird_2013}
A. Bacchelli and C. Bird, ``Expectations, Outcomes, and Challenges of Modern Code Review,'' ICSE, 2013.

\bibitem{mcintosh_2015_impact}
S. McIntosh et al., ``An Empirical Study of the Impact of Modern Code Review Practices on Software Quality,'' EMSE, 2015.

\bibitem{kononenko_2016_review_quality}
O. Kononenko, O. Baysal, and D. M. German, ``Code Review Quality: How Developers See It,'' ICSE, 2016.

\bibitem{dibiase_2016_security_chromium}
M. di Biase, M. Bruntink, and A. Bacchelli, ``A Security Perspective on Code Review: The Case of Chromium,'' SCAM, 2016.

\bibitem{sadowski_2018_google}
C. Sadowski et al., ``Modern Code Review: A Case Study at Google,'' ICSE-SEIP, 2018.

\bibitem{macleod_2018_trenches}
L. MacLeod et al., ``Code Reviewing in the Trenches: Understanding Challenges and Best Practices,'' IEEE Software, 2018.

\bibitem{davila_nunes_2021_slr_taxonomy}
N. Davila and I. Nunes, ``A Systematic Literature Review and Taxonomy of Modern Code Review,'' JSS, 2021.

\bibitem{tufano_2021_towards}
R. Tufano et al., ``Towards Automating Code Review Activities,'' ICSE, 2021.

\bibitem{li_2022_codereviewer}
Z. Li et al., ``Automating Code Review Activities by Large-Scale Pre-training,'' arXiv, 2022.

\bibitem{tufano_2022_pretrained}
R. Tufano et al., ``Using Pre-Trained Models to Boost Code Review Automation,'' arXiv, 2022.

\bibitem{yu_2023_security_defect}
J. Yu et al., ``Security Defect Detection via Code Review,'' arXiv, 2023.

\bibitem{watanabe_2024_chatgpt}
M. Watanabe et al., ``On the Use of ChatGPT for Code Review,'' EASE, 2024.

\bibitem{tang_2024_codeagent}
D. Tang et al., ``CodeAgent: Autonomous Communicative Agents for Code Review,'' EMNLP, 2024.

\bibitem{asthana_2024_whodo}
N. Asthana et al., ``WhoDo: Automating Reviewer Suggestions at Scale,'' ICSE-SEIP, 2024.

\bibitem{frommgen_2024_resolving}
A. Frömmgen et al., ``Resolving Code Review Comments with Machine Learning,'' ICSE-SEIP, 2024.

\bibitem{charoenwet_2024_sast_secure_review}
W. Charoenwet et al., ``An Empirical Study of Static Analysis Tools for Secure Code Review,'' ISSTA, 2024.

\bibitem{olewicki_2024_llm_review_impact}
D. Olewicki et al., ``Impact of LLM-based Review Comment Generation in Practice,'' arXiv, 2024.

\bibitem{pornprasit_2024_gpt35}
C. Pornprasit and C. Tantithamthavorn, ``GPT-3.5 for Code Review Automation,'' arXiv, 2024.

\bibitem{rasheed_2024_ai_powered}
Z. Rasheed et al., ``AI-Powered Code Review with LLMs: Early Results,'' arXiv, 2024.

\bibitem{tufano_2024_strengths_weaknesses}
R. Tufano et al., ``Code Review Automation: Strengths and Weaknesses of the State of the Art,'' arXiv, 2024.

\bibitem{vijayvergiya_2024_autocommenter}
M. Vijayvergiya et al., ``AI-Assisted Assessment of Coding Practices in Modern Code Review,'' arXiv, 2024.

\bibitem{peng_2025_icodereviewer}
Y. Peng et al., ``iCodeReviewer: Improving Secure Code Review with Mixture of Prompts,'' ASE, 2025.

\bibitem{cihan_2025_llm_review_eval}
U. Cihan et al., ``Evaluating Large Language Models for Code Review,'' arXiv, 2025.

\bibitem{hong_2025_rag_review}
H. Hong and Baik, ``Retrieval-Augmented Code Review Comment Generation,'' arXiv, 2025.

\bibitem{liu_2025_noisy_data}
C. Liu et al., ``Too Noisy To Learn: Enhancing Data Quality for Code Review Comment Generation,'' arXiv, 2025.

\bibitem{yang_2026_roadmap}
Z. Yang et al., ``A Roadmap for Modern Code Review: Challenges and Opportunities,'' arXiv, 2026.

% === 漏洞检测与修复新增 ===

\bibitem{kim2017vuddy}
S. Kim et al., ``VUDDY: A Scalable Approach for Vulnerable Code Clone Discovery,'' IEEE S\&P, 2017.

\bibitem{li2018vuldeepecker}
Z. Li et al., ``VulDeePecker: A Deep Learning-Based System for Vulnerability Detection,'' NDSS, 2018.

\bibitem{zhou2019devign}
Y. Zhou et al., ``Devign: Effective Vulnerability Identification by Learning Comprehensive Program Semantics via Graph Neural Networks,'' NeurIPS, 2019.

\bibitem{feng2020codebert}
Z. Feng et al., ``CodeBERT: A Pre-Trained Model for Programming and Natural Languages,'' EMNLP Findings, 2020.

\bibitem{fu2021linevul}
M. Fu and C. Tantithamthavorn, ``LineVul: A Transformer-based Line-Level Vulnerability Prediction,'' MSR, 2021.

\bibitem{ahmad2021plbart}
W.U. Ahmad et al., ``Unified Pre-training for Program Understanding and Generation,'' NAACL, 2021.

\bibitem{ding2023vulhawk}
R. Ding et al., ``VulHawk: Cross-Architecture Vulnerability Detection with Entropy-Based Binary Code Search,'' NDSS, 2023.

\bibitem{peng2023ptlvd}
H. Peng et al., ``PTLVD: Program Slicing and Transformer-based Line-level Vulnerability Detection,'' SCAM, 2023.

\bibitem{mirsky2023vulchecker}
Y. Mirsky et al., ``VulChecker: Graph-based Vulnerability Localization in Source Code,'' USENIX Security, 2023.

\bibitem{woo2023v1scan}
S. Woo et al., ``V1SCAN: Discovering 1-day Vulnerabilities in Reused Open-Source Components,'' USENIX Security, 2023.

\bibitem{chen2024vullibgen}
Y. Chen and X. Xie, ``VulLibGen: Generating Names of Vulnerability-Affected Packages via LLM,'' ACL, 2024.

\bibitem{liu2024contrastive}
T. Liu et al., ``Applying Contrastive Learning to Code Vulnerability Type Classification,'' EMNLP, 2024.

\bibitem{paquette2024llmscannot}
R. Paquette et al., ``LLMs Cannot Reliably Identify and Reason About Security Vulnerabilities,'' IEEE S\&P, 2024.

\bibitem{wang2024taskfusion}
Y. Wang et al., ``Suitable is the Best: Task-Oriented Knowledge Fusion for Vulnerability Detection,'' NeurIPS, 2024.

\bibitem{liu2024eatvul}
P. Liu et al., ``EaTVul: ChatGPT-based Evasion Attack against Software Vulnerability Detection,'' USENIX Security, 2024.

\bibitem{risse2024limits}
N. Risse and M. Böhme, ``Uncovering the Limits of Machine Learning for Automatic Vulnerability Detection,'' USENIX Security, 2024.

\bibitem{shao2024fvddpm}
Z. Shao et al., ``FVD-DPM: Fine-Grained Vulnerability Detection via Diffusion Probabilistic Model,'' USENIX Security, 2024.

\bibitem{shimmi2024vulsim}
A. Shimmi et al., ``VulSim: Multi-Dimensional Neighbor Embeddings for Vulnerability Detection,'' USENIX Security, 2024.

\bibitem{zhang2024vulllmeval}
H. Zhang et al., ``VulnLLMEval: A Framework for Evaluating LLMs in Vulnerability Detection and Patching,'' arXiv, 2024.

\bibitem{zhang2024prompt}
C. Zhang et al., ``Prompt-Enhanced Software Vulnerability Detection Using ChatGPT,'' arXiv, 2024.

\bibitem{zhou2025vulndetrepair}
X. Zhou et al., ``Vulnerability Detection and Repair,'' ACM TOSEM, 2025.

\bibitem{khare2025llmvulndet}
S. Khare et al., ``LLM Vulnerability Detection Effectiveness,'' ICST, 2025.

\bibitem{niu2025largetomammoth}
C. Niu et al., ``From Large to Mammoth: A Comparative Evaluation of LLMs for Vulnerability Detection,'' NDSS, 2025.

\bibitem{song2025nodemedic}
J. Song et al., ``NodeMedic-FINE: Automatic Detection and Exploit Synthesis for Node.js Vulnerabilities,'' NDSS, 2025.

\bibitem{rossi2025crossdomain}
L. Rossi et al., ``Cross-Domain Evaluation of Transformer Models for Vulnerability Detection,'' Springer, 2025.

\bibitem{zhang2025tplpatch}
Y. Zhang et al., ``Enhancing Security of Third-Party Library Reuse: 1-day Vulnerability Code Patching,'' NDSS, 2025.

\bibitem{wen2025appatch}
Y. Wen et al., ``AppAtch: Automated Adaptive Prompting for LLM-based Vulnerability Repair,'' USENIX Security, 2025.

\bibitem{kim2025san2patch}
J. Kim et al., ``Logs In, Patches Out: Automated Vulnerability Repair via Tree-of-Thought with LLM,'' USENIX Security, 2025.

\bibitem{yu2025patchagent}
Z. Yu et al., ``PatchAgent: An LLM-based Autonomous Agent for Vulnerability Repair,'' USENIX Security, 2025.

\bibitem{sok2025avr}
H. Zhang et al., ``Automated Vulnerability Repair: Methods, Tools, and Assessments (SoK),'' USENIX SoK, 2025.

\bibitem{sheng2026llmsecurity}
Z. Sheng et al., ``LLMs in Software Security: A Survey of Vulnerability Detection Techniques and Insights,'' ACM CSUR, 2026.

\bibitem{chen2026lmcode}
Y. Chen et al., ``Language Models for Code Security: A Large-Scale Empirical Study,'' ACM TOSEM, 2026.

\bibitem{kaniewski2026detecting}
M. Kaniewski et al., ``Detecting Software Vulnerabilities with LLMs,'' ACM TOSEM, 2026.

% === Agent 安全 / 提示注入 / 越狱新增 ===

\bibitem{zhan2024injecagent}
Q. Zhan et al., ``InjecAgent: Benchmarking Indirect Prompt Injections in Tool-Integrated LLM Agents,'' ACL Findings, 2024.

\bibitem{liu2024promptinjbench}
Y. Liu et al., ``Formalizing and Benchmarking Prompt Injection Attacks and Defenses,'' USENIX Security, 2024.

\bibitem{kong2024injectbench}
N. Kong, ``InjectBench: An Indirect Prompt Injection Benchmarking Framework,'' Virginia Tech, 2024.

\bibitem{li2024instrfollow}
Z. Li et al., ``Evaluating the Instruction-Following Robustness of LLMs to Prompt Injection,'' EMNLP, 2024.

\bibitem{chen2025canipi}
Y. Chen et al., ``Can Indirect Prompt Injection Attacks Be Detected and Removed?'' ACL, 2025.

\bibitem{kim2025manyshot}
S. Kim et al., ``What Really Matters in Many-Shot Attacks?'' ACL, 2025.

\bibitem{evtimov2025wasp}
I. Evtimov et al., ``WASP: Benchmarking Web Agent Security Against Prompt Injection Attacks,'' NeurIPS, 2025.

\bibitem{hines2024spotlighting}
K. Hines et al., ``Defending Against Indirect Prompt Injection Attacks With Spotlighting,'' CAMLIS, 2024.

\bibitem{ayub2024embedding}
M. Ayub and S. Majumdar, ``Embedding-Based Classifiers Can Detect Prompt Injection Attacks,'' CEUR, 2024.

\bibitem{evertz2024whispers}
J. Evertz et al., ``Whispers in the Machine: Confidentiality in LLM-Integrated Systems,'' IEEE S\&P, 2024.

\bibitem{zhou2024rpo}
A. Zhou et al., ``Robust Prompt Optimization for Defending Language Models Against Jailbreaking Attacks,'' NeurIPS, 2024.

\bibitem{chen2025defenseleverage}
Y. Chen et al., ``Defense Against Prompt Injection Attack by Leveraging Attack Techniques,'' ACL, 2025.

\bibitem{jia2025taskshield}
F. Jia et al., ``Task Shield: Enforcing Task Alignment to Defend Against Indirect Prompt Injection,'' ACL, 2025.

\bibitem{wen2025instrdetect}
T. Wen et al., ``Defending against Indirect Prompt Injection by Instruction Detection,'' EMNLP Findings, 2025.

\bibitem{chen2025struq}
S. Chen et al., ``StruQ: Defending Against Prompt Injection with Structured Queries,'' USENIX Security, 2025.

\bibitem{ramakrishnan2025pidefense}
B. Ramakrishnan and A. Balaji, ``Securing AI Agents Against Prompt Injection Attacks,'' arXiv, 2025.

\bibitem{tan2026prarag}
X. Tan et al., ``PRA-RAG: Provably Robust Aggregation in RAG against Retrieval Corruption,'' ACL Findings, 2026.

\bibitem{deng2024masterkey}
G. Deng et al., ``MasterKey: Automated Jailbreaking of Large Language Model Chatbots,'' NDSS, 2024.

\bibitem{yu2024dontlisten}
Z. Yu et al., ``Don't Listen To Me: Understanding and Exploring Jailbreak Prompts of LLMs,'' USENIX Security, 2024.

\bibitem{yu2024fuzzer}
J. Yu et al., ``LLM-Fuzzer: Scaling Assessment of Large Language Model Jailbreaks,'' USENIX Security, 2024.

\bibitem{xu2024bagoftricks}
Z. Xu et al., ``Bag of Tricks: Benchmarking of Jailbreak Attacks on LLMs,'' NeurIPS, 2024.

\bibitem{hu2024adc}
K. Hu et al., ``Efficient LLM Jailbreak via Adaptive Dense-to-Sparse Constrained Optimization,'' NeurIPS, 2024.

\bibitem{zheng2024fewshot}
X. Zheng et al., ``Improved Few-Shot Jailbreaking Can Circumvent Aligned Language Models,'' NeurIPS, 2024.

\bibitem{su2024missionimpossible}
J. Su et al., ``Mission Impossible: A Statistical Perspective on Jailbreaking LLMs,'' NeurIPS, 2024.

\bibitem{gong2025papillon}
X. Gong et al., ``PAPILLON: Efficient and Stealthy Fuzz Testing-Powered Jailbreaks for LLMs,'' USENIX Security, 2025.

\bibitem{krauss2025twinbreak}
T. Krauß et al., ``TwinBreak: Jailbreaking LLM Security Alignments based on Twin Prompts,'' USENIX Security, 2025.

\bibitem{zhang2025tasklevel}
L. Zhang et al., ``Exploiting Task-Level Vulnerabilities: Jailbreak Attack and Defense Benchmarking,'' USENIX Security, 2025.

\bibitem{chao2024jailbreakbench}
P. Chao et al., ``JailbreakBench: An Open Robustness Benchmark for Jailbreaking LLMs,'' NeurIPS, 2024.

\bibitem{han2024wildguard}
S. Han et al., ``WildGuard: Open One-Stop Moderation Tools for Safety Risks, Jailbreaks, and Refusals,'' NeurIPS, 2024.

\bibitem{jiang2024wildteaming}
L. Jiang et al., ``WildTeaming at Scale: From In-the-Wild Jailbreaks to Safer Language Models,'' NeurIPS, 2024.

\bibitem{gong2025safetyMisalign}
Y. Gong et al., ``Safety Misalignment Against Large Language Models,'' NDSS, 2025.

\bibitem{debenedetti2024agentdojo}
E. Debenedetti et al., ``AgentDojo: A Dynamic Environment to Evaluate Prompt Injection Attacks and Defenses for LLM Agents,'' NeurIPS, 2024.

\bibitem{zhang2024asb}
H. Zhang et al., ``Agent Security Bench: Formalizing and Benchmarking Attacks and Defenses in LLM-Based Agents,'' ICLR, 2025.

\bibitem{zhang2026agentaudit}
H. Zhang et al., ``Agent Audit: A Security Analysis System for LLM Agent Applications,'' arXiv, 2026.

\bibitem{wu2025isolateGpt}
Y. Wu et al., ``IsolateGPT: An Execution Isolation Architecture for LLM-Based Agentic Systems,'' NDSS, 2025.

\bibitem{ayzenshteyn2025cloak}
D. Ayzenshteyn et al., ``Cloak, Honey, Trap: Proactive Defenses Against LLM Agents,'' USENIX Security, 2025.

\bibitem{liu2025taintstyle}
F. Liu et al., ``Make Agent Defeat Agent: Taint-Style Vulnerabilities in LLM-Based Agents,'' USENIX Security, 2025.

\bibitem{liu2026malskills}
Y. Liu et al., ``Do Not Mention This to the User: Detecting Malicious Agent Skills,'' USENIX Security, 2026.

\bibitem{wu2025promptleakage}
G. Wu et al., ``Prompt Leakage via KV-Cache Sharing in Multi-Tenant LLM Serving,'' NDSS, 2025.

\bibitem{xu2026duality}
Y. Xu et al., ``LLM Agents Security Duality: A Comprehensive Survey,'' Artificial Intelligence Review, 2026.

\bibitem{khonde2026ragsecurity}
S. Khonde et al., ``End-to-End Security Threats and Defenses in Retrieval-Augmented LLM Agents,'' Discover AI, 2026.

\bibitem{sheng2025cybersecurity_slr}
M. Sheng et al., ``Large Language Models for Cyber Security: A Systematic Literature Review,'' arXiv, 2025.

\bibitem{cruxeval_2024}
``CRUXEval: A Benchmark for Code Reasoning, Understanding and Execution,'' ICML, 2024.

% === 模型完整性与后门 ===

\bibitem{li2023beatrix}
Y. Li et al., ``Beatrix Resurrections: Robust Backdoor Detection via Gram Matrices,'' NDSS, 2023.

\bibitem{pan2023asset}
H. Pan et al., ``ASSET: Robust Backdoor Data Detection across Multiplicity of Deep Learning Paradigms,'' USENIX Security, 2023.

\bibitem{zhang2023pelican}
H. Zhang et al., ``PELICAN: Exploiting Backdoors of Naturally Trained Deep Learning Models in Binary Code Analysis,'' USENIX Security, 2023.

\bibitem{bai2023villain}
J. Bai et al., ``VILLAIN: Backdoor Attacks Against Vertical Split Learning,'' USENIX Security, 2023.

\bibitem{boucher2023trojan}
B. Boucher et al., ``Trojan Source: Invisible Vulnerabilities,'' USENIX Security, 2023.

\bibitem{wei2024lmsanitator}
C. Wei et al., ``LMSanitator: Defending Prompt-Tuning Against Task-Agnostic Backdoors,'' NDSS, 2024.

\bibitem{zhang2024instrbackdoor}
H. Zhang et al., ``Instruction Backdoor Attacks Against Customized LLMs,'' USENIX Security, 2024.

\bibitem{yan2024llmbackdoor}
Y. Yan et al., ``LLM-Assisted Backdoor Attack on Code Completion Models,'' USENIX Security, 2024.

\bibitem{backdooralign2024}
``BackdoorAlign: Mitigating Finetuning-based Jailbreak via Backdoor Safety Alignment,'' NeurIPS, 2024.

\bibitem{philosophers2025}
``Philosophers' Stone: Trojaning Plugins of Large Language Models,'' NDSS, 2025.

\bibitem{barbie2025}
``Barbie: Robust Backdoor Detection via Latent Separability,'' NDSS, 2025.

\bibitem{clibe2025}
``CLIBe: Detecting Dynamic Backdoors in Transformer-based NLP Models,'' NDSS, 2025.

\bibitem{dshield2025}
``DShield: Defending Backdoor Attacks on GNN via Discrepancy Learning,'' NDSS, 2025.

% === LLM 增强安全测试 ===

\bibitem{deng2024pentestgpt}
G. Deng et al., ``PentestGPT: An LLM-empowered Tool for Automated Penetration Testing,'' USENIX Security, 2024.

\bibitem{asmrita2024busybox}
Asmrita et al., ``Fuzzing BusyBox: Leveraging LLM for Crash Reuse in Embedded Bug Unearthing,'' USENIX Security, 2024.

\bibitem{deepgo2024}
``DeepGO: Predictive Directed Greybox Fuzzing,'' NDSS, 2024.

\bibitem{enclavefuzz2024}
``EnclaveFuzz: Finding Vulnerabilities in SGX Applications,'' NDSS, 2024.

\bibitem{llm_protocol_fuzz_2024}
``Large Language Model Guided Protocol Fuzzing,'' NDSS, 2024.

\bibitem{degpt2024}
``DeGPT: Optimizing Decompiler Output with LLM,'' NDSS, 2024.

\bibitem{sigmadiff2024}
``SigmaDiff: Semantics-Aware Deep Graph Matching for Pseudocode Diffing,'' NDSS, 2024.

\bibitem{gnnic2024}
``GNNic: Finding Long Lost Sibling Functions with Abstract Similarity,'' NDSS, 2024.

\bibitem{midas2025}
``The Midas Touch: Triggering Capability of LLMs for RM-API Misuse Detection,'' NDSS, 2025.

\bibitem{api_rules2025}
``Generating API Parameter Security Rules with LLM for API Misuse Detection,'' NDSS, 2025.

\bibitem{yurascanner2025}
``YuraScanner: Leveraging LLMs for Task-Driven Web App Scanning,'' NDSS, 2025.

% === LLM 代码生成安全 ===

\bibitem{pearce2023copilot_correctness}
``Is Your Code Generated by ChatGPT Really Correct? Rigorous Evaluation of LLMs for Code Generation,'' NeurIPS, 2023.

\bibitem{niu2023codexleaks}
C. Niu et al., ``CodexLeaks: Privacy Leaks from Code Generation Language Models in GitHub Copilot,'' USENIX Security, 2023.

\bibitem{indict2024}
``INDICT: Code Generation with Internal Dialogues of Critiques for Both Security and Helpfulness,'' NeurIPS, 2024.

\bibitem{ledex2024}
``LeDex: Training LLMs to Better Self-Debug and Explain Code,'' NeurIPS, 2024.

\bibitem{worldcoder2024}
``WorldCoder: A Model-Based LLM Agent Building World Models by Writing Code,'' NeurIPS, 2024.

\bibitem{verified_transpile2024}
``Verified Code Transpilation with LLMs,'' NeurIPS, 2024.

\bibitem{repocod2025}
``Can Language Models Replace Programmers for Coding? REPOCOD Says 'Not Yet','' ACL, 2025.

\bibitem{trust_copilot2025}
``Can You Really Trust Code Copilot? Evaluating LLMs from Code Security Perspective,'' ACL, 2025.

\bibitem{secrepobench2025}
``SecRepoBench: Benchmarking Code Agents for Secure Code Completion in Real-World Repositories,'' arXiv, 2025.

\bibitem{bars2023fuzztruction}
N. Bars et al., ``Fuzztruction: Using Fault Injection-based Fuzzing to Leverage Implicit Domain Knowledge,'' USENIX Security, 2023.

\bibitem{rajput2023icspatch}
S. Rajput et al., ``ICSPatch: Automated Vulnerability Localization and Non-Intrusive Hotpatching in Industrial Control Systems,'' USENIX Security, 2023.

\bibitem{deng2023nautilus}
Y. Deng et al., ``NAUTILUS: Automated RESTful API Vulnerability Detection with Multi-Step Semantic Graphs,'' USENIX Security, 2023.

\bibitem{deng2023cipherh}
Y. Deng et al., ``CipherH: Automated Detection of Ciphertext Side-Channel Vulnerabilities,'' USENIX Security, 2023.

\bibitem{yuan2023cacheql}
Y. Yuan et al., ``CacheQL: Quantifying and Localizing Cache Side-Channel Vulnerabilities,'' USENIX Security, 2023.

\end{thebibliography}

\end{document}
